\documentclass[11pt, a4paper]{article}

% --- Encodage et langue ---
% Compatible pdflatex / xelatex / lualatex.
% Sous xelatex/lualatex, on charge fontspec (Unicode natif) ; sous pdflatex,
% on garde inputenc/fontenc + lmodern.
\usepackage{iftex}
\ifPDFTeX
  \usepackage[utf8]{inputenc}
  \usepackage[T1]{fontenc}
  \usepackage{lmodern}
  % Caractères Unicode rencontrés dans le mémoire mais non gérés par défaut.
  \DeclareUnicodeCharacter{2248}{\ensuremath{\approx}}
  \DeclareUnicodeCharacter{2265}{\ensuremath{\geq}}
  \DeclareUnicodeCharacter{2264}{\ensuremath{\leq}}
  \DeclareUnicodeCharacter{2260}{\ensuremath{\neq}}
  \DeclareUnicodeCharacter{2212}{\ensuremath{-}}
  \DeclareUnicodeCharacter{2192}{\ensuremath{\rightarrow}}
  \DeclareUnicodeCharacter{2190}{\ensuremath{\leftarrow}}
  \DeclareUnicodeCharacter{2191}{\ensuremath{\uparrow}}
  \DeclareUnicodeCharacter{2193}{\ensuremath{\downarrow}}
  \DeclareUnicodeCharacter{2194}{\ensuremath{\leftrightarrow}}
  \DeclareUnicodeCharacter{21D2}{\ensuremath{\Rightarrow}}
  \DeclareUnicodeCharacter{2026}{\ldots}
  \DeclareUnicodeCharacter{2075}{\ensuremath{^{5}}}
  \DeclareUnicodeCharacter{2076}{\ensuremath{^{6}}}
  \DeclareUnicodeCharacter{25BA}{\ensuremath{\blacktriangleright}}
  \DeclareUnicodeCharacter{25C4}{\ensuremath{\blacktriangleleft}}
  \DeclareUnicodeCharacter{2500}{-}
  \DeclareUnicodeCharacter{2502}{|}
  \DeclareUnicodeCharacter{250C}{+}
  \DeclareUnicodeCharacter{2510}{+}
  \DeclareUnicodeCharacter{2514}{+}
  \DeclareUnicodeCharacter{2518}{+}
  \DeclareUnicodeCharacter{2082}{\textsubscript{2}}
\else
  \usepackage{fontspec}
  \defaultfontfeatures{Ligatures=TeX}
  % « Arial » (équivalent libre métrique-compatible fourni par texlive-fonts-recommended).
  \setmainfont{TeX Gyre Heros}
  \setsansfont{TeX Gyre Heros}
  \setmonofont{TeX Gyre Cursor}
\fi
\usepackage[french]{babel}

% --- Caractères Unicode manquants (valable pdflatex ET xelatex/lualatex) ---
\usepackage{newunicodechar}
\newunicodechar{≈}{\ensuremath{\approx}}
\newunicodechar{≥}{\ensuremath{\geq}}
\newunicodechar{≤}{\ensuremath{\leq}}
\newunicodechar{≠}{\ensuremath{\neq}}
\newunicodechar{⇒}{\ensuremath{\Rightarrow}}
\newunicodechar{→}{\ensuremath{\rightarrow}}
\newunicodechar{←}{\ensuremath{\leftarrow}}
\newunicodechar{↑}{\ensuremath{\uparrow}}
\newunicodechar{↓}{\ensuremath{\downarrow}}
\newunicodechar{↔}{\ensuremath{\leftrightarrow}}
\newunicodechar{−}{\ensuremath{-}}
\newunicodechar{…}{\ldots}
\newunicodechar{⁵}{\ensuremath{^{5}}}
\newunicodechar{⁶}{\ensuremath{^{6}}}
\newunicodechar{►}{\ensuremath{\blacktriangleright}}
\newunicodechar{◄}{\ensuremath{\blacktriangleleft}}
\newunicodechar{─}{\texttt{-}}
\newunicodechar{│}{\texttt{|}}
\newunicodechar{┌}{\texttt{+}}
\newunicodechar{┐}{\texttt{+}}
\newunicodechar{└}{\texttt{+}}
\newunicodechar{┘}{\texttt{+}}
\newunicodechar{₂}{\textsubscript{2}}

% --- Mise en page ---
\usepackage[margin=2.5cm]{geometry}
\usepackage{setspace}
\setstretch{1.15}

% --- Titres : forcer \paragraph et \subparagraph à passer à la ligne ---
% Par défaut LaTeX rend \paragraph (= #### Markdown) en titre « run-in »,
% collé au début du paragraphe qui suit. On les redéfinit en blocs.
\usepackage{titlesec}
\titleformat{\paragraph}[block]
  {\normalfont\normalsize\bfseries}{\theparagraph}{1em}{}
\titlespacing*{\paragraph}
  {0pt}{2.25ex plus 1ex minus .2ex}{1ex plus .2ex}
\titleformat{\subparagraph}[block]
  {\normalfont\normalsize\bfseries\itshape}{\thesubparagraph}{1em}{}
\titlespacing*{\subparagraph}
  {0pt}{2ex plus 1ex minus .2ex}{0.8ex plus .2ex}

% --- Anti-débordement (inline code, longs identifiants, URLs) ---
% Pandoc convertit `code` en \texttt{...}. Par défaut, la fonte monospace
% interdit toute coupure : les identifiants type « recursive-1024-128 »
% débordent en marge.
% Stratégie :
%   (a) marge d'élasticité d'urgence + \sloppy,
%   (b) coupure aux tirets « - » réactivée dans la fonte monospace,
%   (c) le paquet « underscore » rend « \_ » sécable en mode texte,
%       ce qui permet le retour à la ligne après un « _ » dans les
%       identifiants type « entité_émettrice », sans couper le mot
%       au milieu d'une syllabe.
\setlength{\emergencystretch}{4em}
\usepackage[strings]{underscore}

\makeatletter
% Active la coupure aux tirets « - » dans la fonte monospace
\renewcommand\ttfamily{%
  \not@math@alphabet\ttfamily\mathtt
  \fontfamily\ttdefault\selectfont
  \hyphenchar\font=`\-\relax}
\makeatother

% --- Maths ---
\usepackage{amsmath, amssymb, amsfonts}

% --- Tableaux et figures ---
\usepackage{array}
\usepackage{calc}
\usepackage{booktabs}
\usepackage{longtable}
\usepackage{graphicx}
\usepackage{float}

% --- Liens ---
\usepackage[colorlinks=true, linkcolor=blue!60!black, citecolor=blue!60!black, urlcolor=blue!60!black]{hyperref}

% --- Support pour les tables sans légende Pandoc 3.8+ ---
\newcounter{none}
\newcommand{\theHnone}{\thenone}

% --- Code ---
% Build uses pandoc --no-highlight, so we provide minimal Shaded / Highlighting
% environments ourselves (pandoc still wraps code blocks in them).
\usepackage{listings}
\usepackage{xcolor}
\usepackage{fancyvrb}
\usepackage{framed}
\definecolor{shadecolor}{RGB}{248,248,248}
\newenvironment{Shaded}{\begin{snugshade}}{\end{snugshade}}
\DefineVerbatimEnvironment{Highlighting}{Verbatim}{commandchars=\\\{\}}

% --- Notes de bas de page (footnotes Pandoc) ---
\usepackage{footnote}

% --- En-têtes / pieds de page ---
\usepackage{fancyhdr}
\pagestyle{fancy}
\setlength{\headheight}{28pt}
\fancyhf{}
\fancyhead[L]{\nouppercase{\leftmark}}
\fancyhead[R]{\thepage}
\renewcommand{\headrulewidth}{0.4pt}
% Classe « article » : \leftmark est alimenté par \sectionmark.
\renewcommand{\sectionmark}[1]{\markboth{#1}{}}

% --- Titre ---
\title{Jumeaux Numériques et Intelligence Artificielle Topologique :
Vers une prédiction instantanée de la résilience et des émissions
urbaines}
\author{Thomas Clerc}
\date{2026}

% --- Pandoc tight list ---
% Pandoc insère \tightlist au début des listes serrées (items sans ligne
% vide entre eux). La définition par défaut écrase tout espace vertical,
% ce qui rend les puces illisibles. On la neutralise pour conserver les
% espacements standards de \itemize.
\providecommand{\tightlist}{}

% --- Pandoc syntax highlighting macros ---

% --- CSL refs (Pandoc --citeproc) ---

\begin{document}

% --- Page de titre ---
\begin{titlepage}
\centering
\vspace*{2cm}
{\Huge\bfseries Jumeaux Numériques et Intelligence Artificielle
Topologique : Vers une prédiction instantanée de la résilience et des
émissions urbaines\par}
\vspace{1.5cm}
{\Large Mémoire\par}
\vspace{1cm}
{\large Thomas Clerc\par}
\vspace{0.5cm}
{\large 2026\par}
\vfill
\end{titlepage}

% --- Table des matières ---
\tableofcontents
\newpage

% --- Corps ---
\newpage

\section{INTRODUCTION GÉNÉRALE}\label{introduction-guxe9nuxe9rale}

\subsubsection{Contexte Global : Anthropocène, Urbanisation Accélérée et
Électrification des
Mobilités}\label{contexte-global-anthropocuxe8ne-urbanisation-accuxe9luxe9ruxe9e-et-uxe9lectrification-des-mobilituxe9s}

L'humanité est entrée de plain-pied dans l'ère de l'Anthropocène, une
époque géologique inédite caractérisée par l'empreinte prédominante et
irréversible des activités humaines sur l'ensemble des systèmes
géophysiques terrestres. Au cœur de cette grande accélération se trouve
le phénomène urbain. Les villes, bien qu'elles n'occupent
géographiquement que 3 \% de la surface émergée de la Terre, concentrent
aujourd'hui plus de 55 \% de la population mondiale --- un chiffre qui
devrait atteindre 68 \% d'ici 2050. Cette concentration démographique
fait des zones urbaines les principaux moteurs du changement climatique
global : elles consomment près de 78 \% de l'énergie primaire mondiale
et génèrent plus de 70 \% des émissions anthropiques de gaz à effet de
serre. Dans cette équation environnementale critique, le secteur des
transports routiers représente le nœud le plus complexe à délier, étant
responsable à lui seul de près de 25 \% des émissions de dioxyde de
carbone (\(CO_2\)) à l'échelle planétaire.

Parallèlement à cette urgence climatique globale, les dynamiques
contemporaines d'urbanisation se caractérisent par des processus de
densification d'une rapidité et d'une intensité sans précedent,
particulièrement visibles au sein des économies émergentes d'Asie du
Sud-Est. Le Vietnam, engagé dans une transition économique majeure,
illustre de manière paradigmatique cette métamorphose. La capitale,
Hanoï, voit sa périphérie se structurer autour de méga-projets
d'aménagement intégrés appelés ``master-planned communities''. Le
complexe résidentiel et commercial \emph{Vinhomes Ocean Park}, conçu
pour accueillir à terme près de 90 000 résidents sur des surfaces
autrefois à vocation agricole, est le symbole de cette verticalisation
et de cette densification extrêmes. Actuellement habité par environ 60
000 personnes, ce tissu urbain dense concentre des flux de mobilité
massifs et hétérogènes sur des voiries fermées, engendrant une
congestion récurrente et une dégradation accélérée de la qualité de
l'air local.

Pour faire face à la double contrainte de l'explosion des besoins de
mobilité et de l'impératif de décarbonation, la transition vers
l'électrification globale des flottes s'est imposée comme le pivot
central des stratégies publiques et industrielles de transport. Sous
l'impulsion d'acteurs industriels nationaux d'envergure globale comme le
constructeur \emph{VinFast}, le Vietnam s'est engagé dans un plan de
conversion accéléré de ses vecteurs de transport (mobilité individuelle
par deux-roues électriques, véhicules particuliers et bus électriques
pour le transport de masse).

Toutefois, cette transition technologique ultra-rapide se heurte à un
verrou infrastructurel majeur. L'implantation de hubs de recharge à
haute puissance (High Power Charging Hubs), indispensables pour
maintenir la viabilité opérationnelle des flottes électriques en milieu
dense, perturbe profondément le fonctionnement du réseau viaire
existant. La recharge rapide de batteries implique des puissances de
l'ordre de plusieurs mégawatts concentrées sur des points précis du
territoire. L'aménagement urbain moderne doit donc aujourd'hui composer
avec des systèmes hybrides complexes, où les flux cinématiques de
véhicules convergent vers des ressources énergétiques ponctuelles,
créant des interactions non-linéaires entre congestion routière,
contraintes de réseau électrique et dynamique de recharge.

\subsubsection{Le Verrou Technique : Lenteur Computationnelle des
Micro-Simulateurs et Inertie
Décisionnelle}\label{le-verrou-technique-lenteur-computationnelle-des-micro-simulateurs-et-inertie-duxe9cisionnelle}

Dans ce contexte d'hybridation des contraintes urbaines, la
planification et la prise de décision souffrent historiquement d'une
forte inertie. Les outils d'ingénierie du trafic traditionnels,
principalement basés sur des modélisations macroscopiques statiques ou
des approximations fluides simplifiées, s'avèrent incapables de capturer
les comportements hautement non-linéaires et les instabilités
caractéristiques des réseaux saturés. Des phénomènes critiques tels que
les ondes de choc de freinage arrière, les situations de blocage
géométrique complet aux intersections (\emph{gridlocks}), ou la
dynamique fine des files d'attente à l'accès d'une infrastructure de
recharge ne peuvent être appréhendés par des approches globales
moyennes.

Pour répondre à ce besoin de précision, la recherche en ingénierie des
transports a développé des frameworks de micro-simulation dynamique
multi-agents haute-fidélité, au premier rang desquels figure le
progiciel open-source \textbf{SUMO} (Simulation of Urban MObility). Ces
simulateurs microscopiques permettent de modéliser le comportement de
chaque véhicule individuel seconde par seconde, en intégrant des lois
physiques rigoureuses de poursuite de véhicules (car-following) et de
changement de voie. Cette approche granulaire offre une précision
inégalée pour estimer les temps de parcours réels, identifier les
goulots d'étranglement locaux et quantifier les émissions instantanées
de polluants à partir de banques de données de facteurs d'émissions
calibrées (HBEFA3).

Cependant, cette fidélité de modélisation physique se heurte à une
barrière computationnelle majeure. La micro-simulation multi-agent est
un processus intrinsèquement séquentiel et lourd en termes de calcul
processeur. L'évaluation continue des équations différentielles
cinématiques pour chaque agent, le calcul des distances de sécurité, la
résolution des conflits géométriques aux intersections complexes et la
gestion des états de signalisation dynamique représentent une charge de
calcul phénoménale.

Lorsqu'un planificateur urbain doit concevoir l'aménagement optimal
d'une infrastructure critique --- comme l'emplacement et le
dimensionnement d'une station de recharge rapide dans une zone
résidentielle dense --- il fait face à un problème d'optimisation
combinatoire complexe. Tester 1 000 configurations d'aménagement
possibles à l'aide de simulations microscopiques classiques requiert des
centaines d'heures de calcul processeur continu. De plus, à l'échelle de
réseaux métropolitains de grande taille (dont les fichiers géométriques
dépassent fréquemment le gigaoctet), la micro-simulation sature la
mémoire vive des ordinateurs de bureau standard, provoquant des lenteurs
extrêmes liées aux écritures disque virtuelles (mécanisme de SWAP).
Cette contrainte temporelle interdit l'usage des simulateurs physiques
pour la décision interactive ou pour des boucles d'optimisation
automatique en temps réel. Le développement d'outils combinant la
fidélité de la micro-simulation et la réactivité instantanée est le défi
technique majeur traité dans ce mémoire.

\subsubsection{Double Objectif de la Recherche et Positionnement
Académique}\label{double-objectif-de-la-recherche-et-positionnement-acaduxe9mique}

Afin de briser ce verrou technologique sans consentir à des
simplifications dégradant la pertinence scientifique des résultats, ce
mémoire de fin d'étude se structure autour d'un double objectif
méthodologique complémentaire. Nous proposons de dépasser la frontière
traditionnelle entre la précision lente locale et l'approximation rapide
globale en articulant nos recherches selon deux axes distincts et
synergiques.

Le premier axe vise à démontrer l'efficacité et la rigueur de la
micro-simulation microscopique pour l'évaluation chirurgicale
d'infrastructures locales complexes. En partant de réseaux routiers
authentiques modélisés sous forme de graphes topologiques orientés, nous
développons des jumeaux numériques multi-agents dynamiques intégrant des
physiques de déplacement avancées et des modèles de combustion/émissions
validés (HBEFA3). Ce premier volet se concentre sur l'acquisition de
données réelles et la calibration physique fine du comportement
cinématique des véhicules face aux infrastructures de recharge rapide.

Le second axe de recherche propose une rupture conceptuelle majeure. Il
s'agit de s'affranchir définitivement du coût de calcul de la simulation
physique en introduisant le concept d'Intelligence Artificielle
Topologique Spectrale. L'hypothèse scientifique fondatrice de cet axe
stipule que la structure géométrique et topologique intrinsèque du
graphe routier (son squelette mathématique) contient de manière
déterministe les signatures de sa dynamique sous charge (niveaux de
congestion et volume d'émissions de \(CO_2\)). En extrayant les
signatures spectrales de la matrice d'adjacence non-symétrique associée
au réseau urbain et en exploitant des concepts issus de l'analyse
non-normale et de la théorie des perturbations d'opérateurs, nous
construisons un moteur prédictif capable de se substituer au simulateur
physique en une fraction de seconde, offrant un outil d'aide à la
décision d'une agilité inédite pour les planificateurs de la ville
intelligente.

\subsubsection{Annonce Détaillée de la Structure du
Mémoire}\label{annonce-duxe9tailluxe9e-de-la-structure-du-muxe9moire}

Pour exposer cette recherche avec toute la rigueur universitaire
requise, ce mémoire est structuré selon le plan logique suivant :

La \textbf{Première Partie} est dédiée à la collecte et à l'analyse des
données de terrain. Elle décrit le protocole d'observation mis en œuvre
à Vinhomes Ocean Park (Hanoï) pour acquérir la vérité terrain de la
mobilité sous forte densité, présente les technologies de vision par
ordinateur utilisées pour la détection des flottes hétérogènes, et
détaille le processus d'audit de données qui a révélé des biais
systématiques de détection ainsi que les dynamiques de trafic
spécifiques observées lors des périodes de congés.

La \textbf{Deuxième Partie} pose les fondations techniques et théoriques
de nos modèles. D'une part, elle détaille le fonctionnement interne du
simulateur SUMO, notamment les équations cinématiques de poursuite de
véhicules (modèle de Krauß) et les étapes algorithmiques de préparation
des réseaux routiers (calcul des composantes fortement connexes et
filtrage des micro-trajets). D'autre part, elle présente le formalisme
de l'analyse spectrale des graphes asymétriques et de la théorie des
perturbations d'opérateurs qui constituent la signature mathématique
exploitée par notre modèle d'intelligence artificielle prédictive.

La \textbf{Troisième Partie} confronte ces outils théoriques à des
applications concrètes. Elle présente en premier lieu l'étude de cas
localisée du hub de recharge de Sao Bien à Hanoï, modélisant les
dynamiques de files d'attente et le comportement stochastique des
sessions de recharge. En second lieu, elle déploie un cadre expérimental
global pour tester la généricité de notre moteur de simulation sur six
réseaux urbains mondiaux de morphologies contrastées, comparant les
exigences computationnelles (temps CPU, empreinte mémoire RAM) de la
micro-simulation physique face à notre modèle d'intelligence
artificielle spectral.

\newpage

\section{PREMIÈRE PARTIE : ACQUISITION DES DONNÉES, PROTOCOLE DE
COLLECTE ET CONTEXTES DE
MOBILITÉ}\label{premiuxe8re-partie-acquisition-des-donnuxe9es-protocole-de-collecte-et-contextes-de-mobilituxe9}

\subsection{Chapitre 1 : Contexte démographique et modélisation de la
transition
électrique}\label{chapitre-1-contexte-duxe9mographique-et-moduxe9lisation-de-la-transition-uxe9lectrique}

\subsubsection{Le paradigme des villes intelligentes et de la
décarbonation des systèmes de
transport}\label{le-paradigme-des-villes-intelligentes-et-de-la-duxe9carbonation-des-systuxe8mes-de-transport}

La transition vers la ville intelligente (``Smart City'') ne se résume
pas à l'intégration superficielle de technologies de communication au
sein de l'espace urbain ; elle constitue une réponse structurelle à
l'impératif de décarbonation. Dans l'architecture d'une métropole
moderne, les systèmes de transport représentent l'un des principaux
réservoirs d'optimisation carbone. Les politiques publiques
traditionnelles, basées sur l'expansion continue des capacités de
voirie, ont montré leurs limites en se heurtant au phénomène de demande
induite. La décarbonation requiert donc une approche combinée : la
transition technologique des moteurs thermiques vers des systèmes de
propulsion électrique, et l'optimisation dynamique des flux pour
minimiser les pertes énergétiques associées à la congestion.

Cette restructuration des systèmes de transport impose de repenser la
relation entre l'espace routier et le réseau d'alimentation en énergie.
Alors que le ravitaillement en carburant fossile s'effectuait via un
réseau de distribution décentralisé et déconnecté du trafic immédiat
(stations-services hors voirie principale), l'électrification ancre le
ravitaillement au cœur du flux de circulation. Les bornes de recharge
s'insèrent dans le domaine public, sur les places de stationnement ou
les voies de desserte commerciale. Cette hybridation physique transforme
chaque point de recharge en un nœud d'attraction cinématique, modifiant
localement la trajectoire des véhicules et perturbant l'écoulement des
flux environnants.

\subsubsection{Le cadre d'étude vietnamien : Vinhomes Ocean Park (VHOP)
et sa cinétique de
croissance}\label{le-cadre-duxe9tude-vietnamien-vinhomes-ocean-park-vhop-et-sa-cinuxe9tique-de-croissance}

Le développement rapide de l'Asie du Sud-Est s'accompagne d'une
transition urbaine caractérisée par la construction de villes nouvelles
planifiées en périphérie des centres historiques. Au Vietnam, ce modèle
trouve son illustration dans le projet \emph{Vinhomes Ocean Park}
(VHOP), un complexe multifonctionnel de 420 hectares situé à l'est de
Hanoï. Conçu initialement comme une cité satellite résidentielle pour
soulager la pression démographique du district historique de Hoan Kiem,
VHOP a été dimensionné pour accueillir une population nominale de 90 000
habitants.

Cependant, la cinétique d'occupation de VHOP a dépassé les prévisions
initiales. En 2023, le nombre de résidents permanents a franchi le seuil
des 60 000 personnes, et les projections démographiques révisées
indiquent que le site pourrait accueillir à terme plus de 200 000
résidents d'ici 2030 en raison de l'attractivité des infrastructures
scolaires, médicales et commerciales intégrées.

Cette croissance démographique accélérée se traduit par une pression
sans précédent sur les infrastructures routières locales. En 2023, la
population active présente sur site ne représentait que 30 \% à 33 \% de
la capacité finale projetée. L'analyse des flux actuels montre que le
réseau routier interne, bien que moderne, approche déjà de la saturation
lors des heures de pointe. Anticiper le comportement de ce réseau face à
un triplement de la densité de population constitue une problématique
critique pour éviter la paralysie fonctionnelle de la communauté.

\subsubsection{L'intégration des flottes électriques (VinFast) et la
problématique infrastructurelle des hubs de
recharge}\label{lintuxe9gration-des-flottes-uxe9lectriques-vinfast-et-la-probluxe9matique-infrastructurelle-des-hubs-de-recharge}

Au cœur de la stratégie de mobilité de Vinhomes Ocean Park se trouve
l'écosystème électrique porté par le groupe \emph{Vingroup}, à travers
sa filiale automobile \emph{VinFast} et sa filiale de recharge
\emph{V-Green}. L'électrification à VHOP ne relève pas d'une adoption
individuelle lente, mais d'une intégration industrielle planifiée. Les
transports en commun internes sont assurés par des bus électriques
(\emph{VinBus}), et la flotte de taxis opérant sur le site est
constituée en majorité de véhicules électriques de la compagnie
\emph{GSM (Green and Smart Mobility)}.

Pour soutenir cette flotte captive et inciter les résidents à abandonner
le moteur thermique, un réseau de hubs de recharge rapide a été déployé.
Le cas d'étude de Sao Bien se focalise sur la station de recharge
installée en bordure du \emph{Vincom Mega Mall}, l'un des centres
d'attraction majeurs du complexe. Ce hub est équipé de 12 bornes de
recharge ultra-rapides d'une puissance unitaire de 150 kW DC.

L'intégration d'une telle infrastructure génère des contraintes
opérationnelles sévères. D'une part, la demande électrique cumulée de la
station peut atteindre 1,8 MW en période de pointe, imposant des
contraintes de dimensionnement sur le réseau de distribution moyenne
tension local. D'autre part, la géométrie d'accès à la station,
configurée en stationnement en épi à 135 degrés le long d'une voie de
service étroite, crée des conflits d'usage. Les véhicules en attente de
charge partagent la voirie avec la file des taxis en attente de clients
pour le centre commercial, formant des goulots d'étranglement
cinématiques dès que le taux d'occupation des bornes dépasse 75 \%.

\newpage

\subsection{Chapitre 2 : Méthodologie de collecte de données terrain par
vision par
ordinateur}\label{chapitre-2-muxe9thodologie-de-collecte-de-donnuxe9es-terrain-par-vision-par-ordinateur}

\subsubsection{L'infrastructure d'observation : Stratégie du 3ème étage
(Vincom Mega Mall) et minimisation de
l'occlusion}\label{linfrastructure-dobservation-stratuxe9gie-du-3uxe8me-uxe9tage-vincom-mega-mall-et-minimisation-de-locclusion}

La collecte de données précises sur les flux de trafic en milieu urbain
dense nécessite une méthodologie d'observation rigoureuse pour garantir
la qualité des données d'entrée du jumeau numérique. Pour le site de Sao
Bien, nous avons établi un poste d'observation temporaire au 3ème étage
du \emph{Vincom Mega Mall}, au niveau de la zone de restauration faisant
face au corridor de circulation principal et au hub de recharge.

Ce positionnement en hauteur (angle d'observation incliné entre 30 et 45
degrés par rapport à l'horizontale) répond à une contrainte technique
majeure : \textbf{la minimisation de l'occlusion visuelle}.

Dans un contexte de trafic mixte caractéristique du Vietnam, où les flux
comportent une part importante de deux-roues motorisés circulant de
front avec des autobus et des SUV de gabarit important, une capture au
niveau du sol souffre d'un biais de masquage systématique. Les véhicules
volumineux (tels que les bus électriques) occultent totalement les
motocyclistes et les petites voitures situés sur les voies intérieures
ou dans les angles morts géométriques de la caméra.

La perspective plongeante du 3ème étage permet d'obtenir une vue dégagée
``pseudo-orthogonale'' du réseau, garantissant que chaque trajectoire de
véhicule reste visible tout au long du segment d'analyse, éliminant
ainsi les pertes de suivi et améliorant la précision des comptages.

\subsubsection{Contraintes opérationnelles de prise de vue : Choix du
format Portrait vs Paysage face aux obstacles
géométriques}\label{contraintes-opuxe9rationnelles-de-prise-de-vue-choix-du-format-portrait-vs-paysage-face-aux-obstacles-guxe9omuxe9triques}

L'établissement de la ligne de visée depuis le 3ème étage du centre
commercial a imposé des contraintes géométriques strictes liées à
l'architecture du bâtiment. La présence de piliers de soutien en béton,
de montants de fenêtres en aluminium et de vitrines de magasins
obstruait une grande partie du champ visuel horizontal.

L'analyse comparative des formats de capture a révélé les éléments
suivants : * \textbf{Le format paysage (horizontal) :} Bien qu'adapté
pour capturer la longueur du tronçon routier, il intégrait dans le cadre
plusieurs obstacles physiques majeurs qui divisaient la zone d'intérêt
en sous-sections disjointes, empêchant le suivi continu des trajectoires
par l'algorithme de vision par ordinateur. * \textbf{Le format portrait
(vertical) :} En orientant la caméra verticalement, nous avons pu
aligner le champ de vision principal dans l'espace situé entre deux
piliers consécutifs du bâtiment. Cette configuration a permis de cadrer
de manière ininterrompue les trois zones critiques du corridor : la zone
d'approche amont, la zone d'insertion du hub de recharge, et la zone de
sortie vers le carrefour aval.

Le choix du format portrait a ainsi éliminé les pertes de suivi
dynamique causées par les occultations architecturales fixes, assurant
une continuité des données de trajectoire indispensable pour la
calibration des temps de transit.

\subsubsection{Le pipeline de détection automatique et de classification
de véhicules via
YOLOv8}\label{le-pipeline-de-duxe9tection-automatique-et-de-classification-de-vuxe9hicules-via-yolov8}

Pour automatiser l'extraction des données de trafic à partir des
séquences vidéo haute définition capturées sur site, nous avons déployé
un pipeline de traitement d'images basé sur le modèle de réseau de
neurones convolutifs \textbf{YOLOv8} (You Only Look Once, version 8).

Ce pipeline de vision par ordinateur fonctionne de la manière suivante :
1. \textbf{Segmentation temporelle :} Les vidéos brutes sont découpées
en segments homogènes de 10 minutes pour correspondre aux intervalles
standards d'analyse de trafic. 2. \textbf{Inférence et détection :} Le
modèle YOLOv8, pré-entraîné sur le jeu de données COCO (Common Objects
in Context) et optimisé pour la détection en temps réel, traite les
trames vidéo avec un seuil de confiance de détection fixé à \(0.50\).
Pour chaque trame, l'algorithme génère des boîtes de délimitation
(bounding boxes) autour des objets détectés. 3. \textbf{Classification
catégorielle :} Les objets détectés sont classés selon trois classes de
véhicules pertinentes pour l'étude de la mobilité locale : *
\emph{Standard Car} (véhicules particuliers équipés de moteurs
thermiques). * \emph{Electric Vehicle} (véhicules électriques,
identifiables par leur signature visuelle spécifique et leur plaque
d'immatriculation bleue réservée aux VE au Vietnam). * \emph{Motorcycle}
(deux-roues motorisés). 4. \textbf{Suivi de trajectoire (Tracking) :}
L'algorithme associe un identifiant unique à chaque boîte de
délimitation à l'aide d'un filtre de Kalman et d'une matrice de coût
basée sur le recouvrement spatial des boîtes (Intersection over Union -
IoU) entre trames successives. Cela permet de compter le nombre exact de
véhicules franchissant des lignes virtuelles de comptage définies sur
chaque voie de circulation.

\newpage

\subsection{Chapitre 3 : Traitement, audit et correction des données de
trafic
brutes}\label{chapitre-3-traitement-audit-et-correction-des-donnuxe9es-de-trafic-brutes}

\subsubsection{Analyse des biais de classification de l'IA et
application du facteur de correction
systématique}\label{analyse-des-biais-de-classification-de-lia-et-application-du-facteur-de-correction-systuxe9matique}

Malgré l'efficacité de l'architecture YOLOv8 pour la détection
automatique, la phase d'audit qualité des données a mis en évidence un
biais de classification systématique lors de l'analyse des flux de
véhicules particuliers. En comparant les résultats de l'extraction
automatisée à un comptage manuel de référence effectué sur 3 heures de
vidéo, nous avons identifié une \textbf{surestimation constante de la
classe des voitures de l'ordre de +30 \%}.

L'analyse d'erreur a révélé que ce biais provenait de deux facteurs : *
La confusion visuelle induite par les longs véhicules (SUV familiaux
VinFast de type VF8 et VF9, mini-fourgonnettes et véhicules de livraison
légers) qui, sous certains angles de vue inclinés, étaient fragmentés
par l'algorithme en plusieurs boîtes de délimitation distinctes,
générant de fausses détections multiples pour un unique véhicule
physique. * L'occlusion partielle transitoire causée par le passage de
bus électriques, qui masquaient puis révélaient des véhicules adjacents,
poussant l'algorithme de suivi à réinitialiser l'identifiant du véhicule
et à le comptabiliser une seconde fois.

Pour stabiliser le jeu de données et éliminer ce bruit de détection,
nous avons intégré un \textbf{facteur de correction mathématique
systématique de -30 \%} appliqué à la classe des véhicules particuliers
dans le pipeline de traitement des données. Cette correction a permis
d'aligner les comptages automatisés sur les données physiques réelles
validées manuellement.

\subsubsection{Établissement des baselines de trafic et d'émissions
(Midday vs Rush
Hour)}\label{uxe9tablissement-des-baselines-de-trafic-et-duxe9missions-midday-vs-rush-hour}

L'analyse des données corrigées a permis de définir les caractéristiques
de base de la mobilité sur le corridor de Sao Bien à travers deux
périodes représentatives de la journée : le creux de la mi-journée
(Midday) et le pic de fin d'après-midi (Rush Hour).

\paragraph{Baseline Midday (12:00 PM)}\label{baseline-midday-1200-pm}

Cette période caractérise un régime de trafic régulier et modéré, avec
un débit moyen observé de \textbf{124,4 véhicules par intervalle de 10
minutes}. La composition de la flotte se répartit comme suit : *
\emph{Motorcycles} : 63,6 unités (soit 51,1 \% de la flotte). *
\emph{Electric Vehicles (EV)} : 32,5 unités (soit 26,1 \%). *
\emph{Standard Cars (ICE)} : 28,3 unités (soit 22,8 \%).

Le trafic à midi est fluide, et la vitesse moyenne des véhicules
avoisine les 35 km/h. La demande de recharge au hub Vincom Mega Mall est
stable, maintenant un taux d'occupation des bornes de l'ordre de 75 \%
(9 bornes occupées sur 12).

\paragraph{Baseline Rush Hour (5:00
PM)}\label{baseline-rush-hour-500-pm}

Le pic de fin d'après-midi se caractérise par une augmentation de
\textbf{72 \% du volume global}, atteignant un débit moyen de
\textbf{214,5 véhicules par intervalle de 10 minutes}. La structure de
la flotte se modifie significativement : * \emph{Motorcycles} : 136,4
unités (soit 63,6 \% de la flotte). * \emph{Electric Vehicles (EV)} :
45,5 unités (soit 21,2 \%). * \emph{Standard Cars (ICE)} : 32,6 unités
(soit 15,2 \%).

Le volume élevé de deux-roues motorisés (63,6 \%) sature l'espace de
circulation. Les motos occupent les espaces inter-voies, créant des
trajectoires complexes d'évitement pour les véhicules plus volumineux.
Cette densité ralentit l'accès au hub de recharge, les véhicules
électriques devant manœuvrer à travers un flux dense pour s'insérer dans
les slots en épi.

\subsubsection{Phénoménologie du trafic observée : Le phénomène de
``Holiday Reversal'' et dynamiques de congestion
locales}\label{phuxe9nomuxe9nologie-du-trafic-observuxe9e-le-phuxe9nomuxe8ne-de-holiday-reversal-et-dynamiques-de-congestion-locales}

L'analyse des données de collecte a mis en lumière un comportement de
trafic singulier lors des périodes de congés officiels (notamment les
journées du 30 avril et du 1er mai, correspondant aux vacances
nationales du Jour de la Réunification et de la Fête du Travail au
Vietnam). Nous avons qualifié ce phénomène de \textbf{``Holiday
Reversal''} (inversion de vacances).

Alors que les jours de semaine classiques présentent une domination
nette des deux-roues motorisés (51 \% à 64 \% de la flotte), les jours
fériés révèlent une structure inversée, dite ``Car-Locked''. Les
observations enregistrées à 12h00 lors des vacances indiquent : * Une
baisse significative de la proportion de motos, qui chute à \textbf{34,6
\%} (soit 42,25 unités). * Une augmentation importante de la part des
véhicules électriques, qui grimpe à \textbf{43,3 \%} (soit 53,00
unités). * Une part stable de voitures thermiques à \textbf{19,4 \%}
(soit 23,75 unités).

Ce changement s'explique par deux facteurs d'usage : d'une part, les
résidents privilégient les déplacements en voiture familiale pour se
rendre dans les zones de loisirs et les centres commerciaux de VHOP ;
d'autre part, la flotte de taxis électriques GSM intensifie son activité
pour répondre à la demande des visiteurs externes.

Pour le hub de recharge, cette configuration représente une situation
critique : le taux d'occupation des bornes atteint 100 \% de manière
quasi-continue, et des files d'attente se forment sur la voie de
service, bloquant le couloir d'accès au centre commercial et générant
des ralentissements locaux significatifs.

\newpage

\section{DEUXIÈME PARTIE : FONDATIONS TECHNIQUES ET THÉORIQUES DES
SYSTÈMES DE SIMULATION ET DE LA PRÉDICTION
TOPOLOGIQUE}\label{deuxiuxe8me-partie-fondations-techniques-et-thuxe9oriques-des-systuxe8mes-de-simulation-et-de-la-pruxe9diction-topologique}

\subsection{Chapitre 4 : Le moteur de micro-simulation physique
SUMO}\label{chapitre-4-le-moteur-de-micro-simulation-physique-sumo}

\subsubsection{Abstraction topologique de la voirie (Nœuds, arêtes,
connecteurs
géométriques)}\label{abstraction-topologique-de-la-voirie-nux153uds-aruxeates-connecteurs-guxe9omuxe9triques}

Le progiciel SUMO modélise les réseaux de transport sous forme de
réseaux logiques basés sur la théorie des graphes orientés. Dans ce
formalisme, chaque intersection physique est représentée par un nœud
unique doté d'une géométrie polygonale décrivant sa surface de jonction.
Les tronçons routiers reliant les nœuds sont modélisés par des arêtes,
subdivisées en une ou plusieurs voies de circulation (\emph{lanes}).

Chaque voie possède des attributs géométriques et comportementaux
stricts : * Une polyligne tridimensionnelle décrivant son axe central. *
Une largeur constante (généralement fixée à 3,2 mètres pour les voies
urbaines standards). * Une liste de classes de véhicules autorisées (ex:
\texttt{passenger}, \texttt{taxi}, \texttt{motorcycle}, \texttt{bus},
\texttt{delivery}). * Une vitesse limite supérieure déterminant la
vitesse de référence des agents.

La transition entre deux arêtes consécutives s'effectue via des
\textbf{connecteurs géométriques} (\emph{connections}) définis à
l'intérieur des nœuds. Ces connecteurs lient précisément une voie de
l'arête d'approche à une voie de l'arête de sortie. Ils supportent les
règles de priorité (ex: céder le passage, priorité absolue) et les
configurations de signalisation dynamique (phases de feux), permettant
au simulateur de calculer les trajectoires de croisement sans collision
physique.

\subsubsection{Modèle de poursuite cinématique de Krauß (Vitesse
sécuritaire, temps de réaction, bruit
stochastique)}\label{moduxe8le-de-poursuite-cinuxe9matique-de-krauuxdf-vitesse-suxe9curitaire-temps-de-ruxe9action-bruit-stochastique}

Pour reproduire le mouvement individuel des véhicules le long des
arêtes, SUMO implémente par défaut le modèle comportemental de poursuite
de véhicule développé par \textbf{Krauß}. Ce modèle cinématique calcule
à chaque pas de temps la vitesse optimale d'un véhicule suiveur pour
éviter toute collision avec le véhicule leader, même si ce dernier
décélère brutalement.

Soit un véhicule suiveur \(F\) caractérisé par sa position \(x(t)\) et
sa vitesse \(v(t)\), circulant derrière un véhicule leader \(L\)
caractérisé par sa position \(x_L(t)\) et sa vitesse \(v_L(t)\).
L'intervalle spatial libre (ou gap) séparant les deux véhicules est
défini par : \[g(t) = x_L(t) - x(t) - l_L\] où \(l_L\) est la longueur
physique du véhicule leader.

Le modèle détermine la vitesse sécuritaire \(v_{safe}(t)\) par la
relation suivante :
\[v_{safe}(t) = v_L(t) + \frac{g(t) - v_L(t)\tau}{\frac{v(t) + v_L(t)}{2d} + \tau}\]
où \(\tau\) représente le temps de réaction du conducteur, et \(d\) sa
capacité de décélération maximale.

La vitesse théorique souhaitée pour le pas de temps suivant,
\(v_{target}(t)\), est le minimum parmi les contraintes physiques et
légales :
\[v_{target}(t) = \min\left( V_{max}, v(t) + a \cdot \Delta t, v_{safe}(t) \right)\]
où \(V_{max}\) est la vitesse limite de la voie, et \(a\) est la
capacité d'accélération maximale du véhicule.

Enfin, pour introduire la variabilité des comportements humains (retards
de réaction légers, imprécisions de contrôle), une perturbation
stochastique négative \(\eta\) est soustraite de la vitesse cible pour
obtenir la vitesse finale appliquée à l'agent :
\[v(t + \Delta t) = \max\left( 0, v_{target}(t) - \eta \right)\] où
\(\eta\) est une variable aléatoire distribuée uniformément sur
l'intervalle \([0, \sigma \cdot a \cdot \Delta t]\), le paramètre
\(\sigma \in [0, 1]\) caractérisant le degré d'inattention ou
d'imperfection du conducteur.

\subsubsection{Pipeline de traitement topologique : Compilation des
réseaux OSM, nettoyage par
netconvert}\label{pipeline-de-traitement-topologique-compilation-des-ruxe9seaux-osm-nettoyage-par-netconvert}

La conversion de données topologiques brutes issues d'OpenStreetMap
(OSM) en réseaux de simulation exploitables par SUMO nécessite
l'exécution d'un pipeline de traitement rigoureux via l'outil
\texttt{netconvert}.

Ce pipeline effectue les opérations de nettoyage suivantes : 1.
\textbf{Uniformisation géométrique :} Re-projection des coordonnées
géographiques sphériques (WGS84) vers un système de coordonnées
cartésiennes planes UTM (Universal Transverse Mercator), permettant
d'effectuer les calculs de distance et de vitesse en mètres et mètres
par seconde. 2. \textbf{Simplification des nœuds :} Fusion des grappes
d'intersections complexes et des micro-carrefours d'OSM en nœuds
logiques uniques. Cette opération élimine les segments de voirie
artificiellement courts (inférieurs à 5 mètres) qui s'avèrent
impossibles à gérer pour les modèles cinématiques de car-following. 3.
\textbf{Nettoyage des connexions redondantes :} Suppression des voies
d'accès et des bretelles d'insertion non connectées ou mal définies dans
la base de données source, évitant la création de trajectoires
conflictuelles anormales.

\subsubsection{Connexité réseau : Algorithme de Tarjan pour l'extraction
des Composantes Fortement Connexes
(SCC)}\label{connexituxe9-ruxe9seau-algorithme-de-tarjan-pour-lextraction-des-composantes-fortement-connexes-scc}

L'intégrité topologique d'un réseau routier de simulation dépend de sa
connexité globale. En raison de la présence de règles de circulation
directionnelles (sens uniques) et de potentielles erreurs de
numérisation dans les données géospatiales d'entrée, un graphe routier
orienté brut comporte souvent des composantes déconnectées du flux
principal.

Pour écarter ce problème, notre pipeline met en œuvre
l'\textbf{algorithme des Composantes Fortement Connexes (SCC)} de
Tarjan.

Soit un graphe orienté \(G = (V, E)\). Une composante fortement connexe
de \(G\) est un sous-graphe maximal \(G' = (V', E')\) dans lequel il
existe un chemin orienté reliant chaque paire de nœuds de manière
bidirectionnelle :
\[\forall (u, v) \in V'^2, \quad u \rightsquigarrow v \quad \text{et} \quad v \rightsquigarrow u\]

L'algorithme de Tarjan utilise un parcours en profondeur (DFS) pour
identifier les cycles et partitionner le graphe en composantes fortement
connexes en une seule passe, avec une complexité temporelle optimale :
\[\mathcal{O}(|V| + |E|)\]

Une fois la partition \(\mathcal{C} = \{G_1, G_2, \dots, G_k\}\)
obtenue, le pipeline extrait la composante principale possédant le plus
grand cardinal de nœuds :
\[G_{SCC_{max}} = \arg\max_{G_i \in \mathcal{C}} |V_i|\]

Tous les nœuds et arêtes n'appartenant pas à cette composante principale
sont supprimés du réseau de simulation final. Ce traitement garantit que
chaque véhicule injecté sur le réseau dispose d'un chemin physique
valide pour atteindre n'importe quel point de destination, éliminant les
blocages d'agents et les échecs de routage dynamique.

\subsubsection{Routage dynamique : Filtre de distance minimale pour
l'élimination des micro-trajets
parasitaires}\label{routage-dynamique-filtre-de-distance-minimale-pour-luxe9limination-des-micro-trajets-parasitaires}

Lors de la phase de génération automatique de la demande de trafic
(synthèse des trajets), les points d'origine et de destination sont
distribués aléatoirement sur le graphe épuré. Pour éviter l'apparition
de micro-trajets (déplacements de moins de 300 mètres reliant des
intersections adjacentes), nous implémentons une contrainte de distance
minimale lors de la phase de routage.

Soit un couple origine-destination \((o, d) \in V^2\) sélectionné pour
générer le trajet d'un agent. Le trajet n'est validé et écrit dans le
fichier final de demande que s'il satisfait la condition suivante :
\[D_{Euclidienne}(o, d) = \sqrt{(x_d - x_o)^2 + (y_d - y_o)^2} \ge 300 \text{ mètres}\]

Cette contrainte force le planificateur d'itinéraires à rejeter les
trajets de très courte distance. En éliminant ces mouvements
parasitaires qui se limitent à des phases d'insertion-sortie immédiates,
le filtre garantit que l'ensemble de la flotte simulée s'insère dans les
flux de transit principaux du réseau, interagissant de manière
représentative avec la topologie globale de la voirie et la
signalisation.

\newpage

\subsection{Chapitre 5 : Théorie mathématique du trafic et de la
topologie
spectrale}\label{chapitre-5-thuxe9orie-mathuxe9matique-du-trafic-et-de-la-topologie-spectrale}

L'évaluation de la résilience et des émissions de carbone au sein d'un
réseau viaire repose traditionnellement sur des modèles dynamiques
microscopiques ou sur des approximations macroscopiques fluides. Bien
que ces méthodes offrent une précision locale indéniable, elles
souffrent d'une lenteur computationnelle rédhibitoire et n'offrent que
peu d'indications théoriques a priori sur la vulnérabilité intrinsèque
de la structure. Dans ce chapitre, nous développons un formalisme
mathématique rigoureux qui extrait la ``signature urbaine'' d'une ville
à travers l'analyse spectrale et non-normale de sa matrice d'adjacence.
Ce cadre théorique permet de quantifier l'instabilité potentielle du
trafic et d'identifier les fragilités topologiques structurelles avant
même toute simulation cinématique.

\subsubsection{Formalisation matricielle du réseau
viaire}\label{formalisation-matricielle-du-ruxe9seau-viaire}

Pour modéliser mathématiquement le réseau routier, nous le représentons
sous la forme d'un graphe orienté et pondéré \(G = (V, E)\), où : *
\(V = \{v_1, v_2, \dots, v_n\}\) désigne l'ensemble des nœuds
(\(n = |V|\)), représentant les intersections physiques du réseau. *
\(E \subset V \times V\) désigne l'ensemble des arêtes orientées
(\(m = |E|\)), modélisant les tronçons routiers à sens unique reliant
ces intersections.

La connectivité et l'impédance physique du réseau sont codées dans la
\textbf{matrice d'adjacence pondérée} \(A \in \mathbb{R}^{n \times n}\).
Contrairement aux graphes non orientés simples où la matrice d'adjacence
est symétrique et binaire, la représentation réaliste d'un réseau urbain
impose une double complexité : 1. \textbf{Asymétrie structurelle :}
L'existence de sens uniques et de priorités de passage implique que
l'existence d'un arc \(v_i \to v_j\) n'entraîne pas celle de l'arc
\(v_j \to v_i\). Ainsi, \(A_{ij} \neq A_{ji}\). 2. \textbf{Pondération
physique :} Chaque coefficient \(A_{ij}\) n'indique pas une simple
connexion, mais est défini comme une fonction de la capacité physique du
tronçon routier. Pour un tronçon reliant le nœud \(i\) au nœud \(j\), la
pondération est donnée par :
\[A_{ij} = \begin{cases} \frac{L_{ij} \cdot C_{ij}}{W_{ij}} & \text{si } (v_i, v_j) \in E \\ 0 & \text{sinon} \end{cases}\]
où \(L_{ij}\) représente la longueur du tronçon (en mètres), \(C_{ij}\)
le nombre de voies de circulation, et \(W_{ij}\) la vitesse maximale
autorisée (en m/s). Cette formulation permet d'intégrer les goulets
d'étranglement géométriques directement dans les coefficients de la
matrice.

Pour illustrer ce formalisme, considérons un exemple minimal de réseau
routier à 4 nœuds représentant une intersection en boucle fermée avec un
goulot d'étranglement :

\begin{verbatim}
       (v1) ---------> (v2)
        ^               |
        |               | (Goulot d'étranglement)
        |               v
       (v4) <--------- (v3)
\end{verbatim}

Dans ce modèle simplifié, la matrice d'adjacence orientée non pondérée
\(A_{ex}\) s'écrit : \[A_{ex} = \begin{pmatrix}
0 & 1 & 0 & 0 \\
0 & 0 & 1 & 0 \\
0 & 0 & 0 & 0 \\
1 & 0 & 1 & 0
\end{pmatrix}\] L'absence de symétrie de cette matrice est évidente
(\(A_{ex} \neq A_{ex}^T\)), ce qui introduit des propriétés spectrales
spécifiques détaillées ci-après.

\subsubsection{Le concept de non-normalité et amplification
transitoire}\label{le-concept-de-non-normalituxe9-et-amplification-transitoire}

Un opérateur linéaire ou une matrice carrée \(A\) est dit
\textbf{normal} si et seulement s'il commute avec son adjoint, soit
\(A A^T = A^T A\). Dans le cas des réseaux routiers réels orientés,
cette relation n'est jamais vérifiée : la matrice d'adjacence \(A\) est
intrinsèquement \textbf{non-normale} (\(A A^T \neq A^T A\)).

Cette non-normalité a des implications physiques profondes sur la
dynamique du trafic. Pour une matrice normale, les vecteurs propres
forment une base orthonormale de l'espace, et le comportement à court
terme du système est entièrement dicté par ses valeurs propres. Si
toutes les valeurs propres ont une partie réelle négative (ou un module
inférieur à 1 en discret), toute perturbation décroît de manière
monotone.

Dans un système non-normal, les vecteurs propres ne sont plus
orthogonaux et peuvent devenir presque colinéaires. Par conséquent, même
si toutes les valeurs propres indiquent une stabilité asymptotique à
long terme (\(\rho(A) < 1\)), le système peut subir une
\textbf{amplification transitoire massive} de l'état avant sa
décroissance finale. Ce phénomène est connu sous le nom d'\textbf{effet
coup de bélier} (\emph{water hammer effect}) du trafic : * Une
perturbation locale mineure (comme un retard de 10 secondes à un feu ou
un ralentissement ponctuel) excite les modes non-orthogonaux du réseau.
* Ces modes interfèrent de manière constructive à court terme,
provoquant une amplification géométrique de la file d'attente. * Cette
accumulation transitoire sature les intersections en amont par effet
domino, paralysant temporairement une large zone urbaine et générant un
pic d'émissions de \(CO_2\) avant que le système ne parvienne à évacuer
le flux.

\subsubsection{\texorpdfstring{Le Rayon Spectral (\(\rho\)) et le
Théorème de
Perron-Frobenius}{Le Rayon Spectral (\textbackslash rho) et le Théorème de Perron-Frobenius}}\label{le-rayon-spectral-rho-et-le-thuxe9oruxe8me-de-perron-frobenius}

Le spectre d'une matrice d'adjacence \(A\), noté \(\sigma(A)\), est
l'ensemble de ses valeurs propres \(\lambda_i \in \mathbb{C}\) vérifiant
l'équation caractéristique \(\det(\lambda I - A) = 0\). Le \textbf{rayon
spectral} \(\rho(A)\) est défini comme le module de la plus grande
valeur propre : \[\rho(A) = \max_{\lambda \in \sigma(A)} |\lambda|\]

Puisque les coefficients \(A_{ij}\) de notre matrice d'adjacence
pondérée sont strictement non-négatifs (\(A_{ij} \ge 0\)), nous pouvons
appliquer le \textbf{théorème de Perron-Frobenius}. Si le graphe \(G\)
est fortement connexe (ce qui est garanti par notre filtrage
algorithmique de Tarjan), la matrice \(A\) est irréductible. Le théorème
stipule alors que : 1. Le rayon spectral \(\rho(A)\) est lui-même une
valeur propre de \(A\), simple et strictement positive
(\(\rho(A) > 0\)). 2. Il existe un vecteur propre à droite \(v_{PF}\)
associé à \(\rho(A)\) dont toutes les composantes sont strictement
positives (\(v_{PF} > 0\)), appelé vecteur de Perron-Frobenius. 3. Cette
valeur propre domine toutes les autres :
\(\forall \lambda \in \sigma(A) \setminus \{\rho(A)\}, \ |\lambda| \le \rho(A)\).

Physiquement, le rayon spectral \(\rho(A)\) quantifie la capacité
intrinsèque de transit du réseau. Un rayon spectral élevé indique la
dominance d'un mode principal de circulation, caractéristique des villes
fortement hiérarchisées dotées d'autoroutes urbaines ou de grands
boulevards périphériques (par exemple, Paris ou Madrid). À l'inverse, un
rayon spectral proche de 1 caractérise des réseaux hautement réguliers
et distribués (type grille comme Berlin ou Los Angeles), où les flux
n'ont pas d'axe de concentration obligatoire.

\subsubsection{\texorpdfstring{La Constante de Kreiss (\(K\)) et la
dynamique de
crise}{La Constante de Kreiss (K) et la dynamique de crise}}\label{la-constante-de-kreiss-k-et-la-dynamique-de-crise}

Pour quantifier rigoureusement la sensibilité d'un réseau non-normal aux
amplifications transitoires et modéliser son instabilité dynamique, nous
introduisons la \textbf{constante de Kreiss} \(K(A)\). Soit \(A\) une
matrice stable (\(\rho(A) < 1\)). La constante de Kreiss est définie
mathématiquement par la borne supérieure de la norme de son opérateur
résolvant sur le disque unité :
\[K(A) = \sup_{|z| > 1} (|z| - 1) \left\| (zI - A)^{-1} \right\|\] où
\(\|\cdot\|\) désigne la norme matricielle induite (généralement la
norme spectrale \(L_2\)).

Le théorème des matrices de Kreiss établit des bornes strictes reliant
cette constante à l'amplification transitoire maximale de la puissance
de la matrice, décrivant l'évolution temporelle discrète d'une
perturbation :
\[K(A) \le \sup_{k \ge 0} \left\| A^k \right\| \le e \cdot n \cdot K(A)\]
où \(n\) est la dimension de la matrice.

En ingénierie du trafic, la constante de Kreiss agit comme un
\textbf{détecteur de nervosité} ou de fragilité structurelle du réseau :
* Si \(K(A) \approx 1\), le réseau est robuste et se comporte comme un
système normal. Les perturbations de trafic s'amortissent de manière
monotone sans amplification. * Si \(K(A) \gg 10\), le réseau présente
une vulnérabilité critique. Une infime fluctuation de la demande ou un
ralentissement local va induire une congestion transitoire
disproportionnée. La pollution en \(CO_2\), directement liée aux phases
d'arrêt et de réaccélération des véhicules lors des embouteillages,
devient extrêmement instable et subit des pics de pollution d'autant
plus élevés que la constante de Kreiss est forte.

\subsubsection{\texorpdfstring{Les Normes de Transfert \(H_2\) et
\(H_\infty\)}{Les Normes de Transfert H\_2 et H\_\textbackslash infty}}\label{les-normes-de-transfert-h_2-et-h_infty}

Afin de caractériser la réponse globale du réseau à des scénarios
d'injection continue de trafic, nous modélisons la propagation de la
congestion comme un système dynamique linéaire entrée-sortie. Soit la
fonction de transfert du réseau \(T(z) = (zI - A)^{-1}\). Nous évaluons
cette dynamique via deux normes d'espaces de Hardy :

\paragraph{\texorpdfstring{La Norme \(H_\infty\) (Gain
Maximal)}{La Norme H\_\textbackslash infty (Gain Maximal)}}\label{la-norme-h_infty-gain-maximal}

La norme \(H_\infty\) représente la borne supérieure du gain d'énergie
du système sur l'ensemble des fréquences complexes. Elle correspond au
rayon du plus grand cercle de la résolvante et se définit par :
\[\|T\|_{H_\infty} = \sup_{|z| > 1} \left\| (zI - A)^{-1} \right\|_2 = \sup_{\theta \in [0, 2\pi]} \sigma_{max}\left( (e^{i\theta}I - A)^{-1} \right)\]
où \(\sigma_{max}\) est la valeur singulière maximale.

Physiquement, \(\|T\|_{H_\infty}\) quantifie le \textbf{pire scénario de
congestion possible} que la structure géométrique du réseau peut générer
sous une charge harmonique. Elle représente le niveau d'émissions de
\(CO_2\) ``de base'' et inévitable inhérent à la simple géométrie des
rues de la ville.

\paragraph{\texorpdfstring{La Norme \(H_2\) (Énergie
Totale)}{La Norme H\_2 (Énergie Totale)}}\label{la-norme-h_2-uxe9nergie-totale}

La norme \(H_2\) mesure l'énergie totale dissipée par le système en
réponse à une perturbation impulsionnelle unitaire (par exemple, un
afflux soudain de véhicules à une intersection). Elle est définie par :
\[\|T\|_{H_2} = \left( \frac{1}{2\pi} \int_{0}^{2\pi} \text{Tr}\left( T(e^{i\theta}) T(e^{i\theta})^H \right) d\theta \right)^{1/2} = \left( \sum_{k=0}^{\infty} \left\| A^k \right\|_F^2 \right)^{1/2}\]
où \(\|\cdot\|_F\) désigne la norme de Frobenius.

La norme \(H_2\) reflète la \textbf{mémoire temporelle du réseau}. Une
valeur \(H_2\) élevée indique que la congestion mettra un temps
considérable à s'évacuer après la fin de la perturbation originelle (par
exemple, une persistance anormalement longue des embouteillages et des
émissions associées bien après l'heure de pointe).

\subsubsection{Théorie des perturbations de
Kato}\label{thuxe9orie-des-perturbations-de-kato}

Pour agir sur la robustesse du réseau et concevoir des politiques
d'aménagement optimales (sens uniques, fermetures de voies, feux
dynamiques), il est nécessaire de quantifier l'impact d'une modification
structurelle locale sur les valeurs propres de la matrice d'adjacence.
Nous exploitons pour cela la \textbf{théorie des perturbations des
opérateurs linéaires} formalisée par Tosio Kato.

Soit la matrice d'adjacence nominale \(A\) perturbée par une
modification locale \(\delta A\) proportionnelle à un paramètre réel
\(\epsilon\) : \[A(\epsilon) = A + \epsilon B\] où \(B\) est la matrice
de perturbation modélisant la création ou la suppression d'une liaison
routière.

Pour une valeur propre \(\lambda_i\) non dégénérée de la matrice \(A\),
associée au vecteur propre à droite \(v_i\) et au vecteur propre à
gauche \(w_i\) (tels que \(A v_i = \lambda_i v_i\) et
\(w_i^T A = \lambda_i w_i^T\)), le développement perturbatif de
\(\lambda_i(\epsilon)\) s'écrit :
\[\lambda_i(\epsilon) = \lambda_i + \epsilon \lambda_i^{(1)} + \epsilon^2 \lambda_i^{(2)} + \mathcal{O}(\epsilon^3)\]

\paragraph{Dérivée Première des Valeurs
Propres}\label{duxe9rivuxe9e-premiuxe8re-des-valeurs-propres}

La variation au premier ordre est donnée par la formule classique :
\[\lambda_i^{(1)} = \frac{w_i^T B v_i}{w_i^T v_i}\]

Cette dérivée première indique la direction de sensibilité immédiate.
Elle permet de savoir si l'ajout d'une connexion (ou son renforcement)
va stabiliser ou déstabiliser le réseau.

\paragraph{Dérivée Seconde des Valeurs
Propres}\label{duxe9rivuxe9e-seconde-des-valeurs-propres}

Pour capturer les effets de second ordre et les couplages non-linéaires
induits par la perturbation, nous utilisons la formule de Kato faisant
intervenir la résolvante réduite \(S_i\) :
\[\lambda_i^{(2)} = w_i^T B S_i B v_i\] où \(S_i\) désigne l'opérateur
réduit associé à la valeur propre \(\lambda_i\), défini par :
\[S_i = \lim_{z \to \lambda_i} (zI - A)^{-1} (I - P_i)\] avec
\(P_i = \frac{v_i w_i^T}{w_i^T v_i}\) le projecteur spectral sur le
sous-espace propre associé à \(\lambda_i\).

L'accès à cette dérivée seconde est capital car la non-normalité de
\(A\) induit des variations fortement non-linéaires des valeurs propres
sous perturbation.

\paragraph{Formulation d'une loi de contrôle de
robustesse}\label{formulation-dune-loi-de-contruxf4le-de-robustesse}

En utilisant ces dérivées analytiques, nous formulons une \textbf{loi de
contrôle géométrique} pour optimiser la robustesse du réseau routier.
L'objectif est de trouver la perturbation optimale \(B^*\) (sélectionnée
parmi un ensemble de modifications topologiques admissibles
\(\mathcal{B}\)) qui minimise la constante de Kreiss ou le rayon
spectral dominant afin de maximiser la résilience du graphe :
\[B^* = \arg\min_{B \in \mathcal{B}} \rho(A + \epsilon B) \approx \arg\min_{B \in \mathcal{B}} \left( \lambda_{PF} + \epsilon \frac{w_{PF}^T B v_{PF}}{w_{PF}^T v_{PF}} + \epsilon^2 w_{PF}^T B S_{PF} B v_{PF} \right)\]

Cette formulation analytique évite d'avoir à recalculer le spectre
complet pour chaque modification virtuelle, offrant un guide
mathématique pour reconcevoir la signalisation ou le sens de circulation
d'une ville de façon optimale et rationnelle.

\newpage

\subsection{Chapitre 6 : Le moteur d'intelligence artificielle
prédictif}\label{chapitre-6-le-moteur-dintelligence-artificielle-pruxe9dictif}

Pour dépasser le coût computationnel prohibitif des simulateurs
physiques multi-agents comme SUMO, ce chapitre présente la conception et
l'évaluation d'un métamodèle basé sur l'intelligence artificielle. En
exploitant la ``signature urbaine'' spectrale et topologique développée
au Chapitre 5, ce moteur prédictif estime de manière instantanée (de
l'ordre de la milliseconde) les émissions globales de \(CO_2\) sous
charge de trafic, ouvrant la voie à des outils d'aide à la décision en
temps réel pour l'aménagement du territoire.

\subsubsection{Objectif scientifique du
métamodèle}\label{objectif-scientifique-du-muxe9tamoduxe8le}

La micro-simulation dynamique (SUMO) résout les équations cinématiques
de poursuite de véhicules seconde par seconde pour chaque agent. Ce
niveau de détail est nécessaire pour calibrer précisément les
comportements locaux, mais il engendre un coût de calcul exponentiel
avec la taille du réseau et la charge de véhicules. À l'échelle de
réseaux métropolitains ou pour des études d'optimisation nécessitant
l'évaluation de milliers de scénarios, cette lenteur limite fortement
l'interactivité.

L'objectif de notre approche est de construire un métamodèle
d'apprentissage supervisé capable d'apprendre la fonction de transfert
implicite du simulateur physique : \[f : \mathcal{X} \to \mathcal{Y}\]
où : * \(\mathcal{X}\) est l'espace des caractéristiques structurelles
et spectrales d'une ville complété par la charge de trafic appliquée. *
\(\mathcal{Y}\) représente les émissions totales de \(CO_2\) (en
kilogrammes) ou la vitesse moyenne globale (en m/s) générées par la
simulation SUMO de référence (Ground Truth).

Une fois entraîné, ce métamodèle permet de court-circuiter l'étape de
simulation physique en fournissant des prédictions immédiates avec une
perte de précision minime.

\subsubsection{Architecture et mathématiques de
XGBoost}\label{architecture-et-mathuxe9matiques-de-xgboost}

Pour modéliser cette fonction complexe, nous avons sélectionné
l'algorithme de boosting d'arbres de décision régularisé
\textbf{XGBoost} (eXtreme Gradient Boosting). Cet algorithme construit
de manière séquentielle un ensemble d'arbres de régression additifs.

Soit un jeu de données d'entraînement
\(\mathcal{D} = \{(x_i, y_i)\}_{i=1}^N\). À l'étape \(t\), la prédiction
du modèle pour l'échantillon \(i\) s'écrit sous la forme additive :
\[\hat{y}_i^{(t)} = \hat{y}_i^{(t-1)} + f_t(x_i)\] où
\(f_t \in \mathcal{F}\) représente la structure du nouvel arbre de
décision ajouté au modèle à l'étape \(t\).

La fonction de coût objective \(\mathcal{L}^{(t)}\) à minimiser lors de
l'apprentissage de ce \(t\)-ième arbre intègre une pénalité de
régularisation stricte :
\[\mathcal{L}^{(t)} = \sum_{i=1}^N l\left(y_i, \hat{y}_i^{(t-1)} + f_t(x_i)\right) + \Omega(f_t)\]

\paragraph{\texorpdfstring{Terme de régularisation
\(\Omega(f)\)}{Terme de régularisation \textbackslash Omega(f)}}\label{terme-de-ruxe9gularisation-omegaf}

Pour éviter le sur-apprentissage (overfitting), particulièrement sur des
jeux de données de taille modérée, le terme de régularisation pénalise
la complexité de l'arbre \(f\) en fonction de son nombre de feuilles
\(T\) et des poids associés à ces feuilles \(w \in \mathbb{R}^T\)
(régularisation hybride L1/L2) :
\[\Omega(f_t) = \gamma T + \frac{1}{2} \lambda \sum_{j=1}^T w_j^2 + \alpha \sum_{j=1}^T |w_j|\]
où \(\gamma\) contrôle le coût d'ajout d'une nouvelle feuille, tandis
que \(\lambda\) et \(\alpha\) sont les coefficients de régularisation
respectifs L2 (Ridge) et L1 (Lasso) appliqués aux scores des feuilles.

\paragraph{Approximation par développement de Taylor au second
ordre}\label{approximation-par-duxe9veloppement-de-taylor-au-second-ordre}

Pour optimiser rapidement cette fonction de coût non linéaire, XGBoost
applique un développement de Taylor au second ordre autour de la
prédiction précédente \(\hat{y}_i^{(t-1)}\) :
\[\mathcal{L}^{(t)} \approx \sum_{i=1}^N \left[ l\left(y_i, \hat{y}_i^{(t-1)}\right) + g_i f_t(x_i) + \frac{1}{2} h_i f_t^2(x_i) \right] + \gamma T + \frac{1}{2} \lambda \sum_{j=1}^T w_j^2\]
où \(g_i\) et \(h_i\) désignent respectivement les dérivées première
(gradient) et seconde (hessienne) de la fonction de perte par rapport à
la prédiction courante :
\[g_i = \frac{\partial l\left(y_i, \hat{y}_i^{(t-1)}\right)}{\partial \hat{y}_i^{(t-1)}} \quad \text{et} \quad h_i = \frac{\partial^2 l\left(y_i, \hat{y}_i^{(t-1)}\right)}{\partial \left(\hat{y}_i^{(t-1)}\right)^2}\]

En regroupant les termes par feuille \(j\), le problème d'optimisation
se ramène à la résolution analytique des poids optimaux des feuilles
\(w_j^*\), ce qui confère à XGBoost sa rapidité et sa robustesse face
aux données bruitées ou de dimensions restreintes.

\subsubsection{L'espace des descripteurs (Features) et Signature
Urbaine}\label{lespace-des-descripteurs-features-et-signature-urbaine}

Le vecteur d'entrée \(x_i \in \mathcal{X}\) fourni au modèle
d'intelligence artificielle est structuré de manière à encapsuler la
physique globale et la dynamique spectrale du réseau. Il comporte six
descripteurs fondamentaux : 1. \textbf{\texttt{total\_vehicles} (La
contrainte de charge) :} Le nombre total de véhicules injectés dans le
réseau (de 5 000 à 30 000 agents). C'est la variable d'effort externe.
2. \textbf{\texttt{n\_nodes} (Complexité géométrique) :} Le nombre total
de carrefours, traduisant le nombre de points de conflit potentiels. 3.
\textbf{\texttt{n\_edges} (Capacité d'accueil) :} Le nombre de segments
routiers du réseau. 4. \textbf{\texttt{density} (Densité de maillage) :}
Ratio caractérisant la connectivité physique globale et la disponibilité
d'itinéraires alternatifs. 5. \textbf{\texttt{kreiss\_constant}
(Instabilité de Kreiss - \(K\)) :} Variable spectrale caractérisant la
sensibilité aux amplifications transitoires et aux ondes de choc de
trafic (Chapitre 5). 6. \textbf{\texttt{spectral\_radius} (Rayon
Spectral - \(\rho\)) :} Marqueur de l'asymétrie structurelle et de la
hiérarchisation des corridors de transit (Chapitre 5).

L'intégration conjointe de descripteurs géométriques classiques et de
descripteurs spectraux constitue la clé de voûte de notre modèle : les
premiers fixent l'échelle du réseau tandis que les seconds informent
l'IA sur son comportement dynamique en situation critique.

\subsubsection{Feature Importance et
Sensibilité}\label{feature-importance-et-sensibilituxe9}

Pour comprendre comment l'IA prend ses décisions, nous avons extrait
l'importance relative des descripteurs (\emph{Feature Importance}) via
l'algorithme Random Forest :

\begin{verbatim}
               IMPORTANCE DES DESCRIPTEURS DANS LA PRÉDICTION
  ========================================================================
  [total_vehicles]     |============================= 48% ============================|
  [density]            |======== 14% ========|
  [n_edges]            |======= 13% =======|
  [n_nodes]            |======= 13% =======|
  [kreiss_constant]    |==== 7% ====|
  [spectral_radius]    |=== 5% ===|
  ========================================================================
\end{verbatim}

L'analyse de cette distribution révèle une répartition logique : 1.
\textbf{La Charge Spatiale (\texttt{total\_vehicles} - 48\%) :} Reste le
facteur prédominant de dimensionnement de la pollution. 2. \textbf{La
Topologie Statique (\texttt{density}, \texttt{n\_edges},
\texttt{n\_nodes} - 40\% cumulés) :} Définit l'échelle géométrique et le
nombre de voies alternatives offertes au trafic. 3. \textbf{La Signature
Spectrale (\texttt{kreiss\_constant}, \texttt{spectral\_radius} - 12\%
cumulés) :} Bien que ce poids cumulé de 12\% puisse sembler secondaire,
il constitue la contribution scientifique majeure de notre modèle. Ces
12\% agissent comme un \textbf{ajusteur de crise}. Lorsque deux villes
présentent une charge et une taille similaires, ce sont uniquement ces
indicateurs spectraux (Kreiss et rayon spectral) qui permettent à
XGBoost de déterminer l'apparition d'un point de rupture et de moduler
la prédiction d'émissions de \(CO_2\) à la hausse pour refléter la
congestion induite par la non-normalité topologique.

\subsubsection{Espace de transfert topologique : Généralisation
Cross-City}\label{espace-de-transfert-topologique-guxe9nuxe9ralisation-cross-city}

Le défi ultime de ce moteur de prédiction réside dans sa capacité de
généralisation transversale (\emph{cross-city generalization}) : prédire
les émissions d'une nouvelle ville pour laquelle nous ne disposons
d'aucune simulation physique SUMO (par exemple, Nairobi), uniquement à
partir de ses caractéristiques géométriques et spectrales.

Pour ce faire, nous formalisons l'apprentissage par transfert sous la
forme d'une interpolation barycentrique au sein de l'espace des
signatures spectrales. Soit \(\vec{x}_{target}\) le vecteur des
descripteurs spectraux et topologiques normalisés de la ville cible
(Nairobi), extrait instantanément à partir de ses données OpenStreetMap.

Nous exprimons ce vecteur comme une combinaison linéaire convexe des
vecteurs de descripteurs de \(k\) villes déjà simulées et présentes dans
la base d'entraînement :
\[\vec{x}_{target} = \sum_{j=1}^k \alpha_j \vec{x}_j\] sous les
contraintes de fermeture convexe :
\[\sum_{j=1}^k \alpha_j = 1 \quad \text{et} \quad \alpha_j \ge 0 \quad \forall j \in \{1, \dots, k\}\]

Les coefficients de pondération \(\alpha_j\) sont calculés en résolvant
un problème d'optimisation quadratique de distance minimale dans
l'espace des descripteurs normalisés :
\[\vec{\alpha}^* = \arg\min_{\vec{\alpha}} \left\| \vec{x}_{target} - \sum_{j=1}^k \alpha_j \vec{x}_j \right\|_2^2\]

Si le profil spectral de Nairobi se révèle proche de celui de Marseille
(poids \(\alpha_{Marseille} = 0.70\)) et de celui de Berlin (poids
\(\alpha_{Berlin} = 0.30\)), l'intelligence artificielle est en mesure
de prédire la courbe d'émissions de la ville cible en exploitant les
surfaces de décision non linéaires apprises sur ces analogues
morphologiques. Cette méthode de généralisation s'affranchit entièrement
de la simulation physique multi-agent lente, ramenant le temps
d'évaluation d'un nouveau plan urbain de plusieurs heures à quelques
microsecondes.

\newpage

\section{TROISIÈME PARTIE : APPLICATIONS EXPÉRIMENTALES ET VALIDATION
COMPARATIVE}\label{troisiuxe8me-partie-applications-expuxe9rimentales-et-validation-comparative}

\subsection{Chapitre 7 : Étude de cas localisée -- Le hub de recharge
rapide de Sao Bien
(Hanoï)}\label{chapitre-7-uxe9tude-de-cas-localisuxe9e-le-hub-de-recharge-rapide-de-sao-bien-hanouxef}

\subsubsection{Morphologie géométrique et contraintes d'infrastructure
du site
d'étude}\label{morphologie-guxe9omuxe9trique-et-contraintes-dinfrastructure-du-site-duxe9tude}

Le complexe résidentiel et commercial de Vinhomes Ocean Park (VHOP),
situé à la périphérie Est de la métropole de Hanoï, constitue une
opportunité d'étude unique pour analyser la pénétration à grande échelle
de l'électromobilité en milieu urbain hyper-dense. Notre recherche se
focalise spécifiquement sur le corridor de l'avenue Sao Bien, un axe
routier bidirectionnel structurant qui dessert l'entrée principale du
centre commercial \emph{Vincom Mega Mall}.

Le site d'étude concentre deux infrastructures de service contiguës
situées sur une voie latérale de desserte commerciale à faible capacité
(largeur de voie de 3,5 mètres, limitée à deux voies de circulation
unidirectionnelle après insertion depuis l'avenue principale) : 1.
\textbf{Une zone d'attente de taxis (GSM Taxi Waiting Zone) :}
Positionnée immédiatement en amont du centre commercial, disposant de 16
places de stockage organisées en stationnement parallèle le long de la
chaussée. 2. \textbf{Le Hub de recharge ultra-rapide VinFast (V-Green
Super-fast Charging Hub) :} Situé à environ 75 mètres après
l'intersection d'entrée, configuré avec 12 emplacements de recharge en
épi à 135 degrés par rapport au sens de circulation de la voie de
service. Cette géométrie spécifique impose aux véhicules un recul à 90
degrés pour s'extraire de la borne de recharge, créant un blocage
temporaire complet de la voie de circulation adjacente lors de chaque
manœuvre de départ.

\begin{verbatim}
                                  Avenue Sao Bien (Axe Principal)
   ========================================================================================
             ||  Insertion
             \/
   ----------------------------------------------------------------------------------------
   [Voie de Dessers]  ==>  [Taxi Waiting Zone (16 slots)]  ==>  [EV Charging Hub (12 slots)]
   ----------------------------------------------------------------------------------------
\end{verbatim}

La problématique physique du site réside dans la superposition de ces
usages au sein d'une géométrie contrainte. Les taxis électriques GSM en
attente de prise en charge et les véhicules électriques (particuliers et
taxis) convergeant vers les bornes de recharge partagent la même voie
d'accès. De plus, la forte prédominance des motocyclistes dans le flux
de trafic de Hanoï (qui représente entre 50 \% et 64 \% de la
composition globale) engendre un phénomène de ``tissage'' permanent. Les
deux-roues s'insèrent dans les micro-intervalles géométriques laissés
par les voitures en manœuvre, augmentant la friction cinématique et
ralentissant les manœuvres d'insertion et d'extraction du hub de
recharge.

\subsubsection{Modélisation cinématique et calibrage stochastique de la
recharge}\label{moduxe9lisation-cinuxe9matique-et-calibrage-stochastique-de-la-recharge}

\paragraph{Le modèle de temps d'arrêt stochastique (Stochastic Dwell
Time)}\label{le-moduxe8le-de-temps-darruxeat-stochastique-stochastic-dwell-time}

Pour représenter fidèlement l'impact de la station de recharge sur le
trafic, il est scientifiquement inexact de modéliser le temps de
raccordement des véhicules électriques (\emph{dwell time}) par une
valeur fixe ou une moyenne simpliste. Dans la réalité physique, la durée
d'une session de recharge est une variable stochastique complexe. Elle
dépend de la puissance nominale de la borne (150 kW DC), mais également
de la courbe de charge de la batterie régie par le système de gestion
thermique (Battery Management System - BMS), du niveau de charge initial
du véhicule à son arrivée (\(SoC_{in}\)) et du niveau de charge souhaité
par l'usager à son départ (\(SoC_{out}\)).

Selon les spécifications techniques de VinFast et de l'opérateur de
recharge V-Green, le cycle de charge standard pour une batterie de 42 à
80 kWh (modèles VF e34, VF8, VF9) de 10 \% à 70 \% de \(SoC\) s'établit
à \textbf{31 minutes} sous une puissance stabilisée. Au-delà de 70 \%,
la puissance de charge décroît de manière exponentielle pour préserver
l'intégrité électrochimique des cellules.

Pour intégrer cette variabilité au sein du simulateur SUMO, nous avons
développé un \textbf{Modèle de Temps d'Arrêt Stochastique} basé sur une
loi normale tronquée. La durée d'arrêt \(d\) d'un véhicule électrique
ciblé pour une session de recharge suit la loi :
\[d \sim \mathcal{N}(\mu, \sigma^2)\] calibrée selon quatre profils
opérationnels d'encombrement du hub issus de nos observations : *
\textbf{Profil Léger (Light) :} Modélise des recharges d'appoint
rapides.
\[\mu = 1050 \text{ s } (17.5 \text{ min}), \quad \sigma = 225 \text{ s } (3.75 \text{ min})\]
* \textbf{Profil Modéré (Moderate) :} Calibré sur le benchmark standard
VinFast.
\[\mu = 2100 \text{ s } (35 \text{ min}), \quad \sigma = 300 \text{ s } (5 \text{ min})\]
* \textbf{Profil Lourd (Heavy) :} Modélise les sessions de recharge
complètes ou les ralentissements dus au partage de puissance dynamique
(\emph{power-sharing}) entre bornes adjacentes.
\[\mu = 3600 \text{ s } (60 \text{ min}), \quad \sigma = 450 \text{ s } (7.5 \text{ min})\]
* \textbf{Profil Critique (Critical) :} Représente les situations de
saturation sévère et le comportement d'indiscipline des usagers
(``overstaying''), où les véhicules restent connectés ou stationnés
après la fin effective de leur cycle de charge.
\[\mu = 5400 \text{ s } (90 \text{ min}), \quad \sigma = 900 \text{ s } (15 \text{ min})\]

\paragraph{Le mécanisme d'initialisation dynamique par
``Warm-Start''}\label{le-muxe9canisme-dinitialisation-dynamique-par-warm-start}

Les simulateurs de trafic microscopiques démarrent classiquement dans un
état ``vide et froid'' (\emph{cold-start}), où aucun véhicule n'est
présent sur le réseau à l'instant \(t = 0\). Pour une modélisation de
flux continu, cette approche introduit un biais transitoire important :
durant les 15 à 20 premières minutes de la simulation, le hub de
recharge est vide, ce qui fausse les temps d'attente et sous-estime la
congestion d'accès.

Pour résoudre ce verrou de fidélité, nous avons conçu un algorithme
d'initialisation dynamique par \textbf{``Warm-Start'' (démarrage à
chaud)}. À l'instant \(t=0\), avant l'injection du premier véhicule
actif de la simulation, le hub est pré-peuplé stochastiquement avec des
\textbf{véhicules fantômes} (\emph{ghost vehicles}) occupant un nombre
prédéterminé de bornes de recharge.

Le nombre de bornes occupées à l'initialisation est calibré selon le
niveau de charge du scénario : * \emph{Scénarios fluides / creux
(Midday) :} 7 à 9 slots occupés stochastiquement (taux d'occupation
initial de 58 \% à 75 \%). * \emph{Scénarios de pointe (Rush Hour /
Holiday) :} 11 à 12 slots occupés (taux d'occupation initial de 91 \% à
100 \%).

Pour éviter un départ simultané de ces véhicules qui créerait une onde
de choc artificielle, l'algorithme attribue à chaque véhicule fantôme
\(j\) un \textbf{temps d'occupation résiduel \(t_{res, j}\)} tiré
aléatoirement selon une distribution uniforme :
\[t_{res, j} \sim \mathcal{U}(300, 900) \text{ secondes}\]

Ainsi, dès la première seconde de la simulation physique, les bornes de
recharge présentent un niveau d'encombrement réaliste et libèrent leurs
places de manière échelonnée dans le temps, offrant une reproduction
fidèle des dynamiques d'attente et d'insertion vécues par les
conducteurs arrivant sur site.

\subsubsection{Analyse des données réelles et identification des profils
empiriques}\label{analyse-des-donnuxe9es-ruxe9elles-et-identification-des-profils-empiriques}

\paragraph{Constitution du jeu de données et exclusion systématique des
bus}\label{constitution-du-jeu-de-donnuxe9es-et-exclusion-systuxe9matique-des-bus}

Pour calibrer le jumeau numérique, nous exploitons les comptages de
trafic issus de notre campagne de mesures vidéo par vision artificielle
(YOLOv8) sur l'avenue Sao Bien. Le jeu de données consolidé comporte 72
enregistrements unitaires de 10 minutes, couvrant différentes journées
et tranches horaires entre le 14 avril et le 7 mai 2026.

Dans notre formalisation des flux cinématiques et environnementaux,
\textbf{les bus de transport en commun ont été systématiquement exclus
des statistiques de composition et des temps d'attente du hub}. Cette
décision méthodologique s'appuie sur deux fondements scientifiques : 1.
\textbf{Indépendance fonctionnelle :} Les bus électriques circulants
(flotte \emph{VinBus}) opèrent sur des itinéraires fixes et rigides avec
des fréquences déterminées (2.0 à 3.0 passages par intervalle de 10
minutes en moyenne). Ils ne s'insèrent jamais dans le hub de recharge
des véhicules légers, disposant de leurs propres stations de charge
dédiées au dépôt principal. 2. \textbf{Distorsion statistique :}
Intégrer les bus dans les calculs de répartition globale sous-estimerait
l'importance relative de la flotte de voitures particulières et de
taxis, qui constitue 100 \% de la demande physique pesant sur le hub de
recharge de Sao Bien. Les bus font donc l'objet d'un suivi analytique
décorrélé (voir Annexe 3).

\paragraph{Caractérisation des trois profils
empiriques}\label{caractuxe9risation-des-trois-profils-empiriques}

Le traitement statistique des données réelles a permis d'isoler trois
profils de trafic distincts et hautement contrastés.

\subparagraph{Le Profil de Ligne de Base (Regular Midday
Baseline)}\label{le-profil-de-ligne-de-base-regular-midday-baseline}

Ce profil, calculé à partir de 58 observations stables lors des heures
creuses de la mi-journée (11h00 - 13h00) hors jours fériés, caractérise
un écoulement fluide à volume modéré : * \textbf{Volume moyen total
(excluant les bus) :} \textbf{134,10 véhicules par tranche de 10
minutes}. * \textbf{Répartition modale :} * \emph{Motorcycles
(deux-roues) :} \textbf{67,97 unités} (\textbf{50,7 \%}). Le mode
deux-roues reste majoritaire, reflétant la structure classique des
déplacements quotidiens à Hanoï. * \emph{Standard Cars (ICE) :}
\textbf{30,38 unités} (\textbf{22,7 \%}). * \emph{Electric Vehicles (EV)
:} \textbf{35,74 unités} (\textbf{26,7 \%}). Ce taux d'électrification
de près de 27 \% de la flotte totale s'explique par la forte
concentration de taxis GSM sur la zone commerciale.

\subparagraph{Le Profil d'Heure de Pointe (Regular Evening
Peak)}\label{le-profil-dheure-de-pointe-regular-evening-peak}

Ce profil modélise la surpression cinématique observée en fin de journée
(17h00 - 18h00), marquée par le retour des travailleurs et la
convergence vers les pôles de restauration du centre commercial : *
\textbf{Volume moyen total (excluant les bus) :} \textbf{227,67
véhicules par tranche de 10 minutes} (soit une augmentation de
\textbf{70 \%} du flux par rapport à la baseline). * \textbf{Répartition
modale :} * \emph{Motorcycles (deux-roues) :} \textbf{143,06 unités}
(\textbf{62,8 \%}). La part des deux-roues s'accroît significativement,
exacerbant le phénomène de friction latérale. * \emph{Standard Cars
(ICE) :} \textbf{34,44 unités} (\textbf{15,1 \%}). * \emph{Electric
Vehicles (EV) :} \textbf{50,17 unités} (\textbf{22,0 \%}).

\subparagraph{Le Profil de Rupture : Le phénomène de ``Holiday
Reversal''}\label{le-profil-de-rupture-le-phuxe9nomuxe8ne-de-holiday-reversal}

Les observations menées lors des journées nationales chômées du 30 avril
et du 1er mai ont révélé une anomalie comportementale majeure.
Contrairement aux modèles classiques qui prévoient une simple homothétie
des volumes, nous avons observé une \textbf{inversion modale
structurelle complète (Holiday Reversal)}.

Le volume de trafic total global diminue légèrement pour atteindre
\textbf{117,17 véhicules par tranche de 10 minutes} en moyenne (en
raison du départ de la population ouvrière et étudiante de la ville),
mais la composition de la flotte subit une mutation drastique : *
\emph{Motorcycles (deux-roues) :} \textbf{46,50 unités} (\textbf{39,7
\%}). On observe une baisse de plus de 10 points de pourcentage de la
part des motos par rapport au midi normal. * \emph{Standard Cars (ICE)
:} \textbf{23,83 unités} (\textbf{20,3 \%}). * \emph{Electric Vehicles
(EV) :} \textbf{46,83 unités} (\textbf{40,0 \%}).

Lors de cette période festive, la part des véhicules électriques grimpe
à 40 \% de la flotte totale en circulation. Si l'on restreint l'analyse
aux seuls véhicules à quatre roues (ICE + EV), \textbf{le taux de
pénétration des véhicules électriques atteint le niveau record de 66,3
\%}. Ce phénomène s'explique par l'utilisation massive de la flotte de
taxis GSM par les familles vietnamiennes visitant le centre commercial
et les lagunes artificielles de VHOP, combinée à l'usage des SUV
électriques personnels.

Pour le hub de recharge, ce profil représente un cas d'étude critique de
saturation instantanée, les bornes fonctionnant en continu à 100 \% de
leur capacité nominale avec l'apparition de files d'attente bloquantes
sur la voirie.

\begin{verbatim}
       RÉPARTITION DU FLUIDE MIDDAY (REGULIER VS HOLIDAY REVERSAL)
  ========================================================================
  [Midi Régulier]  |====== Moto (50.7%) =====|=== EV (26.7%) ===|== ICE (22.7%) ==|
  
  [Midi Vacances]  |==== Moto (39.7%) ===|====== EV (40.0%) ======|== ICE (20.3%) ==|
  ========================================================================
\end{verbatim}

\subparagraph{La Période de Transition Pré-Vacances (27 - 29
Avril)}\label{la-puxe9riode-de-transition-pruxe9-vacances-27---29-avril}

Les mesures enregistrées durant les trois jours précédant les fêtes (27,
28 et 29 avril) décrivent un régime transitoire de montée en charge. Le
volume moyen s'établit à \textbf{123,83 véhicules par tranche de 10
minutes}, avec une répartition intermédiaire : * \emph{Motorcycles :}
\textbf{58,00 unités} (\textbf{46,8 \%}). * \emph{Standard Cars (ICE) :}
\textbf{32,33 unités} (\textbf{26,1 \%}). * \emph{Electric Vehicles :}
\textbf{33,50 unités} (\textbf{27,1 \%}).

Cette phase de transition montre que l'inversion modale n'est pas un
événement instantané mais un processus graduel de modification des
usages de mobilité, fournissant des indicateurs précurseurs
indispensables pour anticiper la congestion du réseau d'alimentation
électrique.

\subsubsection{Évaluation des scénarios d'évolution (S0 - S5) et aide à
la
décision}\label{uxe9valuation-des-scuxe9narios-duxe9volution-s0---s5-et-aide-uxe0-la-duxe9cision}

Pour projeter l'impact de la croissance de Vinhomes Ocean Park d'ici
2030, nous avons simulé cinq scénarios d'évolution sur notre jumeau
numérique SUMO :

\paragraph{Scénario S0 \& S1 : Baseline et Heure de pointe historiques
(2023)}\label{scuxe9nario-s0-s1-baseline-et-heure-de-pointe-historiques-2023}

Calibrés directement sur les débits réels observés. Le réseau est stable
dans le scénario S0 (temps d'attente moyen au hub inférieur à 45
secondes pour les EVs). Dans le scénario S1, des congestions localisées
apparaissent sur la voie de service en raison des conflits d'insertion
avec la file d'attente des taxis, mais la fluidité globale sur l'avenue
Sao Bien reste préservée (vitesse moyenne de 32 km/h).

\paragraph{Scénario S2 : Le pic
Holiday}\label{scuxe9nario-s2-le-pic-holiday}

Modélise la configuration du ``Holiday Reversal'' (taux d'adoption EV à
75 \% au sein des véhicules de tourisme). La station de recharge est
saturée dès les premières minutes. Les véhicules électriques en attente
de bornes libres forment une file d'attente de 4 à 6 véhicules sur la
voie de service, provoquant une baisse de la vitesse moyenne sur cet axe
à moins de 12 km/h.

\paragraph{Scénarios S3a \& S3b : Horizon 2028 (60 \% d'occupation
urbaine)}\label{scuxe9narios-s3a-s3b-horizon-2028-60-doccupation-urbaine}

Ces scénarios projettent une augmentation globale des volumes de trafic
(x2.0 pour le midi S3a, x3.2 pour l'heure de pointe S3b) et un taux
d'adoption EV de 80 \% à 85 \%. * \textbf{Résultats :} Dans le scénario
S3a, la file d'attente au hub commence à déborder de la voie de service
pour impacter la voie lente de l'avenue Sao Bien. * Dans le scénario de
pointe S3b, le goulot d'étranglement devient structurel : le temps
d'attente moyen pour un véhicule électrique avant d'accéder à une borne
s'élève à \textbf{18,4 minutes}. La congestion sur l'avenue Sao Bien
remonte sur plus de 150 mètres en amont de l'intersection d'entrée.

\paragraph{Scénarios S4a \& S4b : Horizon 2030 (Saturation à 80 \%
d'occupation)}\label{scuxe9narios-s4a-s4b-horizon-2030-saturation-uxe0-80-doccupation}

Ce scénario modélise le stress-test limite avec une multiplication par
4.0 du volume de trafic global d'heure de pointe et une transition à 100
\% vers l'électromobilité pour les véhicules légers. * \textbf{Résultats
(S4b) :} Le système atteint son point de rupture physique complète. La
voie de service est entièrement paralysée par une file d'attente
ininterrompue de taxis et de véhicules électriques. Les manœuvres
d'extraction des bornes (recul à 135 degrés) bloquent les flux durant
plusieurs minutes. Ce blocage local se propage par effet cascade
(gridlock) sur l'intégralité du carrefour principal de l'avenue Sao
Bien. Le temps d'attente moyen au hub s'envole à \textbf{45,2 minutes},
et la vitesse moyenne sur le réseau s'effondre à \textbf{4,2 km/h}
(vitesse équivalente à de la marche à pied), confirmant la paralysie
fonctionnelle de l'infrastructure.

\paragraph{Scénario S5 : La politique de mitigation par expansion
physique}\label{scuxe9nario-s5-la-politique-de-mitigation-par-expansion-physique}

Face au constat de rupture du scénario S4b, nous avons testé l'impact
d'une décision d'aménagement urbain : l'expansion du hub de recharge à
\textbf{20 slots (+8 bornes de 150 kW DC)} et la création d'une voie
d'insertion dédiée pour la zone d'attente des taxis. * \textbf{Résultats
:} L'évaluation sur le jumeau numérique SUMO démontre l'efficacité de
cette politique de mitigation. L'augmentation de la capacité de
traitement résorbe la file d'attente sur la voie de service. Le temps
d'attente moyen pour accéder à une borne de recharge chute de 45,2
minutes à \textbf{6,8 minutes} (-85 \%). Les manœuvres de sortie en épi,
bien que toujours gênantes, ne bloquent plus le flux principal grâce à
la voie d'évitement supplémentaire. La vitesse moyenne globale sur le
corridor Sao Bien remonte à \textbf{21,5 km/h}, rétablissant un niveau
de service (Level of Service - LOS) acceptable pour la communauté.

Ce cas d'étude démontre de manière empirique l'utilité du jumeau
numérique microscopique comme outil d'aide à la décision publique pour
planifier la transition énergétique des réseaux de transport urbain.

\newpage

\subsection{Chapitre 8 : Cadre d'expérimentation global et comparative
des réseaux
urbains}\label{chapitre-8-cadre-dexpuxe9rimentation-global-et-comparative-des-ruxe9seaux-urbains}

\subsubsection{Conception d'un framework de simulation multi-ville
agnostique}\label{conception-dun-framework-de-simulation-multi-ville-agnostique}

Pour valider l'indépendance de notre moteur de simulation et sa capacité
à traiter des réseaux de complexités variées, nous avons déployé un
protocole expérimental global impliquant six réseaux urbains aux
caractéristiques contrastées : * \textbf{Versailles :} Réseau aéré de
taille modérée (2 308 nœuds, 9 095 arêtes) caractérisé par une structure
géométrique planifiée. * \textbf{Paris :} Réseau historique de grande
taille (15 735 nœuds, 62 378 arêtes) présentant une structure radiale
dense. * \textbf{Madrid :} Réseau métropolitain étendu (61 258 nœuds,
214 400 arêtes) mêlant un centre dense et des autoroutes urbaines. *
\textbf{Berlin :} Graphe urbain de grande dimension (61 135 nœuds, 213
900 arêtes) caractérisé par une structure mixte. * \textbf{Los Angeles
:} Réseau autoroutier et urbain de très grande taille (176 905 nœuds,
619 100 arêtes) structuré sous forme de grille orthogonale. *
\textbf{Hanoï :} Réseau géant (473 209 nœuds, 1 656 200 arêtes)
modélisant une mégapole asiatique à forte densité et topologie
irrégulière.

Le framework est conçu pour exécuter la même logique de simulation
(chargement du réseau, génération de trajets avec contrainte de distance
minimale, exécution physique de SUMO, et extraction des métriques
écologiques HBEFA3) de manière unifiée, en adaptant uniquement les
chemins d'accès système aux fichiers de topologie respectifs.

\subsubsection{Analyse comparative des résiliences géométriques : Radial
européen vs Grille orthogonale vs Boucles
fermées}\label{analyse-comparative-des-ruxe9siliences-guxe9omuxe9triques-radial-europuxe9en-vs-grille-orthogonale-vs-boucles-fermuxe9es}

L'évaluation des quatre scénarios comportementaux (Constant, Rush Hour,
Max Jam, Bottleneck) sur ces six réseaux révèle l'influence déterminante
de la géométrie de la voirie sur la résilience globale du trafic.

\paragraph{Morphologie Radiale Européenne (Paris,
Madrid)}\label{morphologie-radiale-europuxe9enne-paris-madrid}

Ces réseaux se caractérisent par une forte convergence des axes
principaux vers des points centraux (structures en étoile). Lors du
scénario \emph{Bottleneck}, ces structures se révèlent vulnérables : la
saturation d'un axe central remonte rapidement le long des voies d'accès
radiales, bloquant les carrefours en amont. En raison de la densité des
nœuds historiques et du manque de voies rapides de dérivation, les flux
de trafic ne disposent pas d'alternatives viables, provoquant un
effondrement de la vitesse moyenne et une hausse rapide des émissions de
\(CO_2\).

\paragraph{Grille Orthogonale Nord-Américaine (Los
Angeles)}\label{grille-orthogonale-nord-amuxe9ricaine-los-angeles}

La structure régulière de Los Angeles présente un comportement
différent. Grâce à la régularité des carrefours et à la redondance des
itinéraires parallèles, le réseau absorbe mieux les surcharges locales
du scénario \emph{Bottleneck}. Les véhicules se répartissent d'eux-mêmes
sur les axes adjacents via l'algorithme de routage dynamique. Cependant,
cette structure est sensible au scénario \emph{Max Jam} : le volume
important d'intersections régulées par des feux de signalisation crée
des files d'attente successives qui saturent les intersections si
l'injection de véhicules est trop massive.

\paragraph{Structures Fermées et Voies de Service (Sao Bien,
Hanoï)}\label{structures-fermuxe9es-et-voies-de-service-sao-bien-hanouxef}

La morphologie de Hanoï et des complexes de type Vinhomes se caractérise
par des axes sinueux, des voies de service étroites et des contraintes
d'accès (U-turns imposés par des terre-pleins centraux). Cette géométrie
présente une faible résilience face aux variations de charge. La moindre
obstruction locale (par exemple, un véhicule électrique manœuvrant pour
entrer dans une borne de recharge) bloque l'une des deux voies de
circulation disponibles, forçant les motos à se faufiler et ralentissant
l'ensemble du flux. Ces réseaux présentent un comportement non-linéaire
prononcé, passant sans transition d'un état fluide à une congestion
complète.

\subsubsection{Limites matérielles de la micro-simulation : Surcharge
mémoire (RAM), traitement XML et phénomène de
SWAP}\label{limites-matuxe9rielles-de-la-micro-simulation-surcharge-muxe9moire-ram-traitement-xml-et-phuxe9nomuxe8ne-de-swap}

Les expérimentations menées sur ces six villes ont mis en évidence les
limites physiques des systèmes de calcul pour exécuter des jumeaux
numériques microscopiques sur de grands réseaux.

Le tableau suivant résume la taille des fichiers de topologie sur le
disque dur et l'occupation mémoire vive (RAM) correspondante mesurée en
Python lors de l'exécution de l'étape de routage (chargement du réseau
via \texttt{sumolib}) :

{\def\LTcaptype{none} % do not increment counter
\begin{longtable}[]{@{}
  >{\raggedright\arraybackslash}p{(\linewidth - 6\tabcolsep) * \real{0.2105}}
  >{\centering\arraybackslash}p{(\linewidth - 6\tabcolsep) * \real{0.2632}}
  >{\centering\arraybackslash}p{(\linewidth - 6\tabcolsep) * \real{0.2632}}
  >{\centering\arraybackslash}p{(\linewidth - 6\tabcolsep) * \real{0.2632}}@{}}
\toprule\noalign{}
\begin{minipage}[b]{\linewidth}\raggedright
Ville
\end{minipage} & \begin{minipage}[b]{\linewidth}\centering
Taille du fichier net.xml (Disque)
\end{minipage} & \begin{minipage}[b]{\linewidth}\centering
Occupation mémoire RAM (Python DOM)
\end{minipage} & \begin{minipage}[b]{\linewidth}\centering
Temps de chargement initial
\end{minipage} \\
\midrule\noalign{}
\endhead
\bottomrule\noalign{}
\endlastfoot
\textbf{Versailles} & 6,6 Mo & \textasciitilde110 Mo & \textless{} 1
s \\
\textbf{Paris} & 52 Mo & \textasciitilde880 Mo & 4 s \\
\textbf{Madrid} & 184 Mo & \textasciitilde3,1 Go & 18 s \\
\textbf{Los Angeles} & 533 Mo & \textasciitilde8,9 Go & 55 s \\
\textbf{Hanoï} & 1 239 Mo & \textasciitilde18,5 Go & 142 s \\
\end{longtable}
}

Pour les réseaux de Versailles, Paris et Madrid, l'occupation mémoire
reste compatible avec les architectures de PC de bureau standards
(généralement équipés de 16 Go de RAM).

Cependant, pour le réseau de Hanoï (1,2 Go de données XML sur disque),
l'empreinte mémoire atteint \textbf{18,5 Go de RAM}. Sur une machine
disposant de 16 Go de mémoire physique, ce chargement provoque une
saturation immédiate. Pour maintenir le processus actif, le système
d'exploitation Windows active le mécanisme de \textbf{pagination
virtuelle (SWAP)}, transférant les pages mémoire excédentaires vers le
disque de stockage (SSD ou HDD).

Ce recours au stockage de masse ralentit considérablement la vitesse de
calcul. Le processeur passe une partie de ses cycles d'horloge en
attente d'entrées/sorties disque, allongeant le temps de génération des
itinéraires de quelques secondes à plusieurs dizaines de minutes. Cette
limite matérielle montre la difficulté de déployer des simulations
physiques en temps réel pour l'aide à la décision à l'échelle
métropolitaine.

\subsubsection{Validation du modèle prédictif : Évaluation des temps de
réponse et précision (RMSE) IA vs
SUMO}\label{validation-du-moduxe8le-pruxe9dictif-uxe9valuation-des-temps-de-ruxe9ponse-et-pruxe9cision-rmse-ia-vs-sumo}

Pour contourner ces limites matérielles, le modèle prédictif basé sur
l'intelligence artificielle (XGBoost spectral) a été évalué en termes de
précision et de vitesse d'exécution face au simulateur physique SUMO de
référence.

\paragraph{Analyse comparative des temps de
calcul}\label{analyse-comparative-des-temps-de-calcul}

La différence de temps d'exécution entre l'approche physique
microscopique et l'inférence par apprentissage machine met en évidence
l'intérêt de la rupture technologique :

\begin{itemize}
\tightlist
\item
  \textbf{SUMO (Physique) :} Pour le réseau de Paris avec une charge de
  50 000 véhicules, l'exécution complète du pipeline (génération des
  routes via Dijkstra, calcul de la cinématique des agents et parsing
  final du fichier d'émissions XML \texttt{tripinfo.xml}) requiert un
  temps moyen de \textbf{45 minutes} de calcul CPU continu.
\item
  \textbf{Modèle Prédictif IA (XGBoost) :} L'inférence à partir des
  caractéristiques spectrales du graphe routier s'exécute en seulement
  \textbf{0,2 seconde} sur le même matériel, offrant un gain de vitesse
  de calcul significatif.
\end{itemize}

\paragraph{\texorpdfstring{Précision des estimations de pollution
(\(CO_2\))}{Précision des estimations de pollution (CO\_2)}}\label{pruxe9cision-des-estimations-de-pollution-co_2}

La précision du modèle prédictif a été quantifiée en calculant l'erreur
quadratique moyenne (Root Mean Squared Error - RMSE) sur un ensemble de
test indépendant composé de 200 scénarios urbains simulés sous
différentes configurations de volumes et de flottes :
\[RMSE = \sqrt{\frac{1}{N_{test}} \sum_{i=1}^{N_{test}} \left( y_{real, i} - y_{pred, i} \right)^2}\]

Les validations empiriques indiquent que le modèle XGBoost basé sur les
caractéristiques spectrales (rayon spectral, constante de Kreiss, etc.)
parvient à prédire le volume total d'émissions de \(CO_2\) avec un écart
moyen inférieur à \textbf{4,5 \%} par rapport aux résultats calculés par
le moteur physique de SUMO.

Cette précision démontre que la structure géométrique et spectrale du
graphe routier permet de capter la dynamique des flux sans avoir à
simuler individuellement chaque agent physique. Le planificateur urbain
dispose ainsi d'un outil d'évaluation instantané de l'empreinte
environnementale des aménagements projetés, combinant la réactivité de
l'IA et la fidélité de la micro-simulation.

\newpage

\section{CONCLUSION GÉNÉRALE ET
PERSPECTIVES}\label{conclusion-guxe9nuxe9rale-et-perspectives}

\subsubsection{Bilan des contributions
scientifiques}\label{bilan-des-contributions-scientifiques}

Ce travail de recherche a permis de jeter les bases d'un double cadre
méthodologique pour la modélisation de la mobilité urbaine décarbonée.

D'une part, nous avons développé des jumeaux numériques microscopiques
de haute-fidélité capables de modéliser l'impact cinématique local de
l'intégration d'infrastructures de recharge pour véhicules électriques.
Ce travail a été appliqué sur l'axe historique Paris-Versailles ainsi
que sur le quartier résidentiel à haute densité de Sao Bien à Hanoï. Il
a permis d'intégrer des comportements stochastiques de charge et de
documenter des dynamiques de trafic spécifiques comme l'inversion de
vacances (\emph{Holiday Reversal}).

D'autre part, nous avons proposé une méthode de rupture basée sur
l'intelligence artificielle topologique spectrale. En exploitant la
structure mathématique de la matrice d'adjacence du réseau urbain (rayon
spectral, constante de Kreiss, perturbations de Kato) et l'algorithme
XGBoost, nous avons démontré qu'il est possible d'estimer instantanément
la pollution d'une ville sans simuler individuellement chaque véhicule,
contournant ainsi les limites computationnelles traditionnelles.

\subsubsection{La boucle d'optimisation hybride (IA-SUMO) comme
perspective
ultime}\label{la-boucle-doptimisation-hybride-ia-sumo-comme-perspective-ultime}

La complémentarité des deux approches développées dans ce mémoire ouvre
la voie à un cadre d'optimisation urbaine hybride (IA-SUMO) alliant la
rapidité de l'apprentissage machine et la précision chirurgicale de la
simulation physique.

Dans ce schéma opérationnel : 1. Le modèle prédictif basé sur l'IA
(XGBoost) est utilisé en amont pour explorer rapidement de larges
espaces de solutions (par exemple, tester des milliers d'implantations
géométriques de voirie ou de localisations de hubs de recharge). L'IA
évalue chaque configuration en une fraction de seconde, éliminant les
scénarios inefficaces et sélectionnant les configurations optimales. 2.
Les scénarios retenus par le modèle prédictif sont ensuite injectés dans
le jumeau numérique microscopique haute-fidélité (SUMO). Cette
simulation physique détaillée permet d'affiner l'analyse locale
(vérification des files d'attente au mètre près, comportement
d'évitement des deux-roues, impact électrique précis).

Cette boucle hybride permet d'optimiser l'aménagement urbain à grande
échelle tout en conservant une précision physique sur les points
sensibles du réseau.

\subsubsection{Valorisation académique et publications
associées}\label{valorisation-acaduxe9mique-et-publications-associuxe9es}

Les résultats de ce travail de recherche font l'objet d'une valorisation
académique à travers la préparation de deux publications scientifiques
de fin d'année :

\begin{enumerate}
\def\labelenumi{\arabic{enumi}.}
\tightlist
\item
  \textbf{Première Publication (co-écrite avec VinUniversity) :}

  \begin{itemize}
  \tightlist
  \item
    \emph{Titre :} ``Microscopic Traffic Flow and Emission Modeling of
    High-Power Electric Vehicle Charging Infrastructure in Hyper-Dense
    Master-Planned Communities: The Case of Sao Bien, Vinhomes Ocean
    Park.''
  \item
    \emph{Contenu :} Présentation du protocole de capture par vision par
    ordinateur (YOLOv8), de l'architecture du jumeau numérique SUMO
    localisé, et de l'évaluation des scénarios de densification à Sao
    Bien.
  \end{itemize}
\item
  \textbf{Deuxième Publication (co-écrite avec l'École Hexagone) :}

  \begin{itemize}
  \tightlist
  \item
    \emph{Titre :} ``Topological Graph-Spectral Machine Learning for
    Real-Time Urban CO2 Emissions Prediction: Applying Kato's
    Perturbation Theory and Kreiss Constants to Non-Normal Urban
    Networks.''
  \item
    \emph{Contenu :} Formalisation de la méthode de prédiction spectrale
    sur graphes non-symétriques, étude de la stabilité via la constante
    de Kreiss, et évaluation des performances du modèle XGBoost par
    rapport au simulateur de référence.
  \end{itemize}
\end{enumerate}

\end{document}
